\section[Quadratic Compositions and Quartics]{Quadratic Compositions and Quartic Hypercones}
\label{sec:two-quadratic}

Starting with OT--FKM splitting quartics, we develop the two-block
normal form, quadratic Jordan compositions, and a common theorem
for orthogonal quartic data. Its realization table guides the quartic
applications. A Gram sextic extends the construction, while
higher-bidegree octonionic pullbacks exhibit transverse instability.

\subsection{OT--FKM splitting and the two-block normal form}

Allowing the two radii in the balanced OT--FKM sections of
\cite{CuiFKM25} to vary independently gives singular hypercones
with a $(2,2,4)$ Laplacian algebra.

Let $x,y\in\R^\ell$. Choose skew-symmetric matrices
$A_1,\ldots,A_{m-1}$ satisfying
\[
 A_iA_j+A_jA_i=-2\delta_{ij}\Id.
\]
Assume
\begin{equation}\label{bh:range}
 2\leq m\leq\ell-2.
\end{equation}
Only parameters for which this Clifford representation exists are
included. In particular, $A_1$ is a complex structure, so $\ell$ is even.
Set
\begin{gather}\label{bh:quartic-data}
 u=\norm{x}^2,\quad v=\norm{y}^2,\quad R=u+v,\quad T=uv,\qquad
 Q=\ip{x}{y}^2+\sum_{i=1}^{m-1}\ip{A_i x}{y}^2,\\
 c=\frac{m-1}{\ell-2},\quad a=1-2c,\quad b=c(1-c),\quad
 F=Q-cT.\label{bh:centered}
\end{gather}
Here $0<c<1$ and $b>0$; write $C=\{F=0\}$.

\begin{remark}[Relation to the original OT--FKM cone]\label{bh:fkm-relation}
In these coordinates, the Cartan--M\"unzner OT--FKM polynomial \cite{FKM81} is
\[
 \Phi=R^2-2(u-v)^2-8Q
      =-R^2+8(1-c)T-8F.
\]
On $u=v$, the equation $F=0$ is equivalent to
$\Phi=(1-2c)R^2$. Without the balance condition it is instead
$\Phi=-R^2+8(1-c)T$, whose normalized right-hand side varies with
$T/R^2$. The codimension-two cone studied in \cite{CuiFKM25}
is exactly the balanced section $C\cap\{u=v\}$, whose only singular
point is the vertex; the full hypercone has the two singular axes
described below.
\end{remark}


\subsubsection{Three generators of degrees \texorpdfstring{$(2,2,g)$}{(2,2,g)}}\label{sec:two-block-extension}

Quadratic Jordan closure \cite[Proposition~37(b)]{MendesRadeschi}
and Euler-operator bigrading \cite[Proposition~19(a)]{MendesRadeschiQuadratic}
give a two-block, separately harmonic normal form for Laplacian
algebras generated by two quadratics and one higher-degree polynomial.
This includes the OT--FKM splitting quartics; higher-degree examples
appear at the section's end.

\begin{proposition}\label{prop:two-block-normal}
Let $\cA=\R[R,D,F]$ be a graded Laplacian algebra with independent
quadratic generators $R,D$ and a homogeneous generator $F$ of degree
$e>2$, with $F\notin\R[R,D]$. After an orthogonal
change of coordinates and a linear change of $R,D$, its quadratic
part is generated by $u=|x|^2$, $v=|y|^2$, where
$x\in\R^p$, $y\in\R^q$. There is a radial polynomial $H(u,v)$
such that $P=F-H$ is bihomogeneous of some bidegree $(a,b)$,
$a+b=e$, and $\Delta_xP=\Delta_yP=0$.

Suppose that $p,q\geq2$ and both bidegrees are positive. Then closure
is equivalent to identities
\begin{align}
 |\nabla_xP|^2&=a^2\bigl(Ku^{a-1}v^b+
               Lu^{a/2-1}v^{b/2}P\bigr),\label{eq:two-block-x}\\
 |\nabla_yP|^2&=b^2\bigl(Ku^av^{b-1}+
               Lu^{a/2}v^{b/2-1}P\bigr),\label{eq:two-block-y}
\end{align}
for constants $K>0$ and $L\in\R$. The terms involving $L$ are
omitted unless $a,b$ are both even. Necessarily
\begin{equation}\label{eq:two-block-dimensions}
 \frac{p-2}{a}=\frac{q-2}{b}.
\end{equation}
In particular, $P$ restricts to an isoparametric function on
$\Sph^{p-1}\times\Sph^{q-1}$.
\end{proposition}

\begin{proof}
Write $D(z)=\langle Bz,z\rangle$, with $B$ symmetric. The product
rule gives $\Gamma(D,D)\in\cA_2=\operatorname{span}\{R,D\}$,
so $B^2$ is a linear combination of $B$ and the identity. Since
$D$ is independent of $R$, $B$ has exactly two eigenvalues.
Its orthogonal eigenspaces give $u,v$.

Put $\mathcal B=\R[u,v]$. The operator
$x\cdot\nabla_x=\Gamma(u,\cdot)/2$ preserves $\cA$, and
$\cA_e=\mathcal B_e+\R F$. Decomposing $F$ by bidegree therefore
shows that all its nonradial terms have the same $x$-degree $a$.
Subtract the other radial terms to obtain a bihomogeneous generator
$P_0$ of bidegree $(a,b)$. This algebra is now bigraded. The distinct
bidegree components of $\Delta P_0$ belong to $\mathcal B$, since
they have degree less than $e$.

They vanish unless $a,b$ are even. In the even mixed case write
$\Delta_xP_0=\lambda_xu^{a/2-1}v^{b/2}$ and
$\Delta_yP_0=\lambda_yu^{a/2}v^{b/2-1}$. Commuting these partial
Laplacians gives
$\lambda_x b(b+q-2)=\lambda_y a(a+p-2)$.
Subtracting the appropriate multiple of $u^{a/2}v^{b/2}$ kills
both. If one bidegree is zero, the subtraction only uses the active
block. This proves the separately harmonic normal form.

The two bidegree components of $\Gamma(P,P)$ also lie in the
algebra. Their degree is $2e-2$, so their expressions contain at
most the first power of $P$. Matching bidegrees gives
\eqref{eq:two-block-x}--\eqref{eq:two-block-y} initially with
possibly different constant and linear coefficients. On the unit
product spheres, let $f=P|_{\Sph^{p-1}\times\Sph^{q-1}}$.
The separate tangential gradient squares are quadratics in $f$,
with leading coefficients $-a^2$ and $-b^2$. Its mean is zero,
by separate harmonicity. Its maximum $M>0$ and minimum $m<0$
are roots of both quadratics, since both partial gradients vanish
at these points. Thus the two quadratics are respectively
$a^2(K+Lf-f^2)$ and $b^2(K+Lf-f^2)$, with
$K=-mM>0$ and $L=m+M$. This proves the common coefficients.
Conversely, the two identities and the separate Euler and Laplace
identities give every generator entry of the closure table, so
\cref{prop:MR-generator-test} proves sufficiency.

If $\mu_2$ is the average of $f^2$, integration by parts on each
sphere gives
\[
 a^2K=a(2a+p-2)\mu_2,\qquad
 b^2K=b(2b+q-2)\mu_2.
\]
This implies \eqref{eq:two-block-dimensions}. Finally,
\[
 |\nabla f|^2=(a^2+b^2)(K+Lf-f^2),\qquad
 \Delta f=-[a(a+p-2)+b(b+q-2)]f,
\]
which are the isoparametric equations on the product.
\end{proof}

The block dimensions $p,q$ need
not agree; if $p=q>2$, \eqref{eq:two-block-dimensions} forces $a=b$.
The change $F\mapsto F-H(u,v)$ preserves the algebra but generally
changes its zero set, so the minimal level must still be determined
within each explicit family. If one bidegree is zero, the active block
reduces to \cite[Proposition~37(c)]{MendesRadeschi}, with an independent
Euclidean factor adjoined.



\subsection{Quadratic Jordan algebras and the composition equation}
\label{sec:quadratic-jordan-framework}

Quadratic-generated Laplacian algebras organize compositions of
arbitrary even degree. For a real symmetric matrix $A$, write
$q_A(x)=\langle Ax,x\rangle$.
Then
\begin{equation}\label{eq:quadratic-jordan-correspondence}
 \Delta q_A=2\tr A,\qquad
 \Gamma(q_A,q_B)=4q_{A\circ B},\qquad
 A\circ B=\tfrac12(AB+BA).
\end{equation}
Thus a quadratic-generated algebra $\cA$, containing $R=|x|^2$,
is Laplacian exactly when its quadratic part $J=\cA_2$ is a unital
Jordan algebra under $f\circ g=\Gamma(f,g)/4$
\cite[Proposition~37(b)]{MendesRadeschi}. The matrix realization matters: its multiplicities enter $\Delta q_A$.
The Jordan classification gives Clifford and real, complex, or
quaternionic Gram models; see
\cite[Theorems~B--C, Section~3 and Appendix~A]{MendesRadeschiQuadratic}.

The next proposition is the low-dimensional consequence of this
classification; it concerns algebras generated entirely in degree two.

\begin{proposition}[Three and four quadratic generators]
\label{prop:three-quadratic-normal-forms}
\label{prop:low-dimensional-quadratic-jordan}
Let $\cA$ be generated by homogeneous quadratics, contain $R$, and
be Laplacian. If $k=\dim\cA_2$ is three or four, the following
list is complete, up to orthogonal source coordinates and a linear
change of quadratic generators.
\begin{center}
\small
\begin{tabular}{cl>{\raggedright\arraybackslash}p{0.54\textwidth}}
\toprule
$k$ & $J$ & Quadratic coordinates \\
\midrule
3 & $\R^3$ & $|x_1|^2,|x_2|^2,|x_3|^2$ on three nonzero orthogonal blocks. \\
3 & $\operatorname{Sym}_2(\R)$ & $|x|^2,|y|^2,\langle x,y\rangle$, with $x,y\in\R^\ell$. \\
4 & $\R^4$ & Four nonzero orthogonal block radii. \\
4 & $\R\oplus\operatorname{Sym}_2(\R)$ & An independent block radius and the preceding real Gram coordinates. \\
4 & $\operatorname{Herm}_2(\mathbb C)$ & $|z|^2,|w|^2,\Re\langle z,w\rangle,\Im\langle z,w\rangle$, with $z,w\in\mathbb C^\ell$. \\
\bottomrule
\end{tabular}
\end{center}
Here $k$ is the number of minimal quadratic algebra generators;
algebraic relations may occur among them.
\end{proposition}

\begin{proof}
The correspondence \eqref{eq:quadratic-jordan-correspondence} embeds
$J$ as a unital Jordan subalgebra of the symmetric matrices. In the
formally real Jordan classification recalled in
\cite[Appendix~A]{MendesRadeschiQuadratic}, the simple factors of
real dimension at most four are $\R$, $\operatorname{Sym}_2(\R)$,
and $\operatorname{Herm}_2(\mathbb C)$, of dimensions $1,3,4$.
The spin factors of dimensions three and four are the latter two
matrix algebras. Taking direct sums gives precisely the displayed
list. A primitive one-dimensional factor acts by an orthogonal
projection; these projections split the source. The representation
argument in \cite[Section~3 and Appendix~A]{MendesRadeschiQuadratic}
gives repeated standard real or complex modules for the remaining
factors, hence their Gram coordinates. Finally, a degree-two element
of a quadratic-generated algebra is a linear combination of its
quadratic generators, proving the assertion about their minimal number.
\end{proof}

For the real two-column model use
$R=|x|^2+|y|^2$, $p=|x|^2-|y|^2$, and $q=2\langle x,y\rangle$.
Its image is $\{R\geq0:p^2+q^2\leq R^2\}$ for $\ell\geq2$;
for $\ell=1$ only the boundary $R^2=p^2+q^2$ is attained.
The complex two-column model similarly has coordinates
$R,p,q_1,q_2$, with $q_1+iq_2=2\langle z,w\rangle$ and
$p=|z|^2-|w|^2$. Its image is
$\{R\geq0:p^2+q_1^2+q_2^2\leq R^2\}$ for $\ell\geq2$;
for $\ell=1$ only its boundary is attained. These statements follow by realizing positive
semidefinite $2\times2$ Gram matrices of the allowed rank.

An arbitrary quadratic span need not be closed: on $\R^3$,
$\operatorname{span}\{R,x_1^2,x_1x_2\}$ omits
$\Gamma(x_1x_2,x_1x_2)=x_1^2+x_2^2$. Adjoining $x_2^2$ gives
$\R\oplus\operatorname{Sym}_2(\R)$. Every finite list has a unital
matrix Jordan closure of dimension at most $N(N+1)/2$; restricting
its size is essential, since all even-degree homogeneous polynomials
are polynomials in the full collection $x_ix_j$.

Choose a basis $z=(z_1,\ldots,z_k)$ of $J$ and put
\begin{equation}\label{eq:quadratic-jordan-profile-data}
 \begin{gathered}
 B_{ij}(z)=\tfrac14\Gamma(z_i,z_j),\qquad
 b_i=\tfrac14\Delta z_i,\\
 A_P=\sum_{i,j}B_{ij}P_iP_j,\qquad
 \mathscr D P=\sum_{i,j}B_{ij}P_{ij}+\sum_i b_iP_i.
 \end{gathered}
\end{equation}
The entries of $B$ are linear and the $b_i$ are constant.
The quotient calculus in \cref{thm:finite-reduction} gives
\begin{equation}\label{eq:quadratic-jordan-profile-pullback}
 |\nabla P(z)|^2=4A_P(z),\qquad
 \Delta P(z)=4(\mathscr D P)(z),\qquad
 \cL(P(z))=8\mathscr T(P)(z),
\end{equation}
where
\begin{equation}\label{eq:three-quadratic-profile-numerator}
 \mathscr T(P)=2A_P\mathscr D P
       -\sum_{i,j}B_{ij}P_i\partial_jA_P.
\end{equation}
These identities are evaluated on the actual generator image, also
when the generators satisfy relations.

\begin{proposition}[Composition criterion]
\label{prop:three-quadratic-profile-criterion}
\label{prop:quadratic-jordan-profile-criterion}
Let $P$ be a real homogeneous polynomial of degree $d$ and set
$F=P(z)$. Its regular zero set is minimal if and only if
$\mathscr T(P)=0$ on the part of the generator image where
$P=0$ and $A_P>0$. In particular, $P\mid\mathscr T(P)$ is
sufficient. It is also necessary if $P$ is irreducible and has a
regular zero in an open subset of the image where $B$ is positive
definite. In that case the quotient has degree $2d-2$, while $F$
has degree $2d$.
\end{proposition}

\begin{proof}
The first assertion follows from
\eqref{eq:quadratic-jordan-profile-pullback} and the level-set
mean-curvature identity. Positive definiteness makes the generator
map a submersion and gives $A_P>0$ at regular profile zeros.
Minimality therefore makes $\mathscr T(P)$ vanish on a relatively
open smooth real patch of $P=0$. The patch-divisibility argument in
\cref{prop:minimizing-irreducibility} gives $P\mid\mathscr T(P)$.
Its degree is $3d-2$, which proves the last assertion.
\end{proof}

For orthogonal radii $u_i=|x_i|^2$, with $n_i=\dim V_i$, the data are
\begin{equation}\label{eq:three-radii-profile-operator}
 B=\operatorname{diag}(u_i),\qquad
 \mathscr D P=\sum_i u_iP_{ii}+\tfrac12\sum_i n_iP_i.
\end{equation}
Thus the criterion specializes to
\begin{equation}\label{eq:three-radii-minimality}
 P\ \bigm|\
 \left(\sum_i u_iP_i^2\right)
 \left(2\sum_i u_iP_{ii}+\sum_i n_iP_i\right)
 -\sum_i u_iP_i\partial_i\left(\sum_j u_jP_j^2\right),
\end{equation}
under the irreducibility and interior-zero hypotheses above.

\subsubsection{Clifford profiles with a radial variable}

For a symmetric Clifford system on $\R^N$ with $s\geq2$, let
$p_i(x)=\langle P_ix,x\rangle$ for $1\leq i\leq s$.
The algebra $\R[R,p_1,\ldots,p_s]$ has
\begin{equation}\label{eq:three-spin-profile-operator}
 B=\begin{pmatrix}R&p^T\\p&R\Id_s\end{pmatrix},\qquad
 \mathscr D P=R(P_{RR}+\Delta_pP)
      +2\sum_i p_iP_{Rp_i}+\tfrac N2P_R.
\end{equation}
This follows from \eqref{eq:Clifford-pullback-identities}; the real
and complex two-column models have $(s,N)=(2,2\ell)$ and $(3,4\ell)$.

If a homogeneous profile $f(p)$ is independent of $R$, the identity
\eqref{eq:pullback-defect} reads
\begin{equation}\label{eq:quadratic-Clifford-profile-transfer}
 \cL_x(f(p))=16R^2(\cL_pf)(p)
             -8d f(p)|\nabla_pf|^2(p),\qquad d=\deg f.
\end{equation}
Thus a degree-$d$ minimal target cone yields a degree-$2d$ minimal
source equation, and a harmonic target gives a harmonic source.
This is the Clifford specialization of Riedler's cone-pullback
theorem \cite[Theorem~5.1.4]{RiedlerThesis25}; its subcalibration
transfer is \cref{thm:clifford-area-transfer}. In particular, minimal
cubic targets supply sextic equations. Effective factors and the
variational conclusions are determined by the corresponding results
in \cref{sec:quartic-application}.

Allowing the radial coordinate also includes the OT--FKM splitting
in \eqref{bh:centered}. For a homogeneous profile $P(t,y)$ of degree
$d$, define
\[
 K_P=|\nabla_yP|^2-P_t^2,\qquad
 \square P=\Delta_yP-P_{tt},\qquad
 \Gamma_L(P,K_P)=\nabla_yP\cdot\nabla_yK_P-P_t(K_P)_t.
\]
On $F=P(R,p)=0$, the same profile table gives
\begin{align}
 |\nabla F|^2&=4R K_P,\notag\\
 \cL(F)&=16R^2\left(K_P\square P
                  -\tfrac12\Gamma_L(P,K_P)\right)
             +8(N-2)RK_PP_t.
 \label{eq:radial-Clifford-profile-equation}
\end{align}
For clarity, Euler's identity first gives, on the whole source,
\[
 |\nabla F|^2=4RK_P+8dPP_t,\qquad
 \Delta F=4R\square P+(2N+8d-8)P_t.
\]
Differentiate the first expression before setting $P=0$;
$\Gamma(F,R)=4dF$ and the homogeneity of $K_P$, of degree $2d-2$,
then give \eqref{eq:radial-Clifford-profile-equation}.
For $P=|y|^2-ct^2$, $0<c<1$, this equation selects
$c=2(s-1)/(N-2)$, the familiar minimal isoparametric level.
The splitting quartic instead has profile
$\tfrac14(\sum_{i=1}^m p_i^2+cp_0^2-cR^2)$, while Riedler's
quartic is independent of $R$. Both therefore belong to the same
composition calculus.

\begin{remark}[An elementary radial profile]
\label{sec:three-quadratic-profiles}
In the real two-column model, the radial variable gives quartic minimal levels: for
$\ell>2$ set
\[
 W=p^2+q^2,\qquad c=\frac1{\ell-1},\qquad
 F=W-cR^2,\qquad \Phi=R^2-2W.
\]
Equation \eqref{eq:radial-Clifford-profile-equation} selects this $c$;
$|\nabla\Phi|^2=16R^3$ and $\Delta\Phi=8(\ell-3)R$ identify
$F=0$ as its minimal isoparametric level. Its smallest algebra is
$\R[R,F]$, since $W=F+cR^2$. This minimal level is covered by
\cref{sec:isoparametric-reduction}.
\end{remark}


\subsubsection{Real and complex Gram profiles}

Let $X=[x_1\ \cdots\ x_r]\in\R^{p\times r}$ and set
$h_{ij}=\langle x_i,x_j\rangle$. The independent entries of
$H=X^TX$ span $\operatorname{Sym}_r(\R)$, with $R=\tr H$, and
\begin{equation}\label{eq:real-Gram-differential-table}
 \begin{split}
 \Delta h_{ij}&=2p\delta_{ij},\\
 \Gamma(h_{ij},h_{kl})
 &=\delta_{ik}h_{jl}+\delta_{il}h_{jk}
   +\delta_{jk}h_{il}+\delta_{jl}h_{ik}.
 \end{split}
\end{equation}
Indeed, $\nabla_{x_a}h_{ij}=\delta_{ai}x_j+\delta_{aj}x_i$;
taking inner products and traces proves the table, including the
diagonal entries. For $Z\in\mathbb C^{p\times r}$, use the real
Euclidean metric and write
$h_{ij}=\sum_a\overline{z_{ai}}z_{aj}$, so $H=Z^*Z$.
The complex-bilinear extension of $\Gamma$ satisfies
\begin{equation}\label{eq:complex-Gram-differential-table}
 \Delta h_{ij}=4p\delta_{ij},\qquad
 \Gamma(h_{ij},h_{kl})
       =2(\delta_{jk}h_{il}+\delta_{il}h_{kj}).
\end{equation}
To verify it, use $\Delta=4\sum_{a,b}\partial_{z_{ab}}
\partial_{\bar z_{ab}}$ and
$\Gamma(f,g)=2\sum_{a,b}(f_{z_{ab}}g_{\bar z_{ab}}
+f_{\bar z_{ab}}g_{z_{ab}})$. Taking real and imaginary parts
gives the real table on $\operatorname{Herm}_r(\mathbb C)$.
These are the standard Gram realizations in
\cite[Section~3 and Section~7]{MendesRadeschiQuadratic}.

Their generator images consist of positive semidefinite Gram matrices
of rank at most $p$. For $p\geq r$ the positive definite part is
open and the Gram map is a submersion there: for a prescribed
symmetric or Hermitian increment $A$, the variation
$\dot X=\tfrac12 XH^{-1}A$, or
$\dot Z=\tfrac12 ZH^{-1}A$, has $\dot H=A$.
Consequently \cref{prop:quadratic-jordan-profile-criterion} applies
directly to profiles with a regular zero in this positive definite
part. For $p<r$ the same differential identities apply on the
rank-constrained image. The constructions below select minimal
profiles from these differential tables.

\subsubsection{Orthogonal quartic blocks and their weights}
\label{sec:Clifford-composition-background}

The OT--FKM and Riedler quartics are quadratic profiles in Clifford
coordinates and the radius. The spin-factor realization of
\cref{sec:quadratic-jordan-framework} gives their common quartic table.
The underlying Hurwitz matrix construction and its reduction to
Clifford representations are described in
\cite[Chapter~0, pp.~2--3; Chapter~1, Lemma~1.5 and Proposition~1.7]{Shapiro2000}.

Let $\mathcal C=(P_0,\ldots,P_m)$ be a symmetric Clifford system on
$\R^{2\ell}$, with $m\geq1$, and put
\[
 R=|z|^2,\qquad p_i(z)=\langle P_i z,z\rangle,
 \qquad p=(p_0,\ldots,p_m).
\]
Their differential table is \eqref{eq:Clifford-pullback-identities},
and \cref{lem:Clifford-algebra-pullback} gives $|p|^2\leq R^2$.

In the product $f\circ g=\Gamma(f,g)/4$, the quadratic space
$J=\operatorname{span}\{R,p_0,\ldots,p_m\}$ has multiplication
\begin{equation}\label{eq:Clifford-composition-quadratic-table}
 R\circ R=R,\qquad R\circ p_i=p_i,\qquad
 p_i\circ p_j=\delta_{ij}R.
\end{equation}
Thus $J$ is a spin-factor Jordan algebra. The Clifford relations
give the following table for its quartic block norms.

\begin{lemma}[Orthogonal quartic blocks]
\label{lem:Clifford-block-quartics}
Partition $\{0,\ldots,m\}$ into nonempty sets $I_1,\ldots,I_s$,
and set
\[
 E_j=\sum_{i\in I_j}p_i^2\quad(1\leq j\leq s),\qquad
 E_*=R^2-\sum_{i=0}^m p_i^2.
\]
Then all these polynomials are nonnegative and
\begin{align}
 &E_*+\sum_{j=1}^s E_j=R^2,\qquad
 \Gamma(E_i,E_j)=16RE_i\delta_{ij},
 \label{eq:Clifford-composition-quartic-table}\\
 &\Delta E_j=8|I_j|R,\qquad
 \Delta E_*=8(\ell-m)R.\nonumber
\end{align}
Here the pairing identity includes the index $*$.
\end{lemma}

\begin{proof}
For disjoint blocks the pairing in
\eqref{eq:Clifford-pullback-identities} gives
$\Gamma(E_i,E_j)=0$; for one block it gives
$\Gamma(E_j,E_j)=16RE_j$ and $\Delta E_j=8|I_j|R$.
Also $\Gamma(R^2,E_j)=16RE_j$, so $E_*$ is orthogonal to every
$E_j$.  Expanding its gradient square gives
$\Gamma(E_*,E_*)=16RE_*$.  Finally,
$\Delta R^2=8(\ell+1)R$ proves the asserted Laplacian of $E_*$.
\end{proof}

\begin{remark}[Geometric interpretation]
\label{rem:weighted-Clifford-foliation}
For $\ell>m+1$, the Clifford map on the unit sphere has connected
fibers and hemisphere quotient metric. If all $|I_j|\geq2$, pulling
back the target $\prod_j\mathrm{SO}(|I_j|)$-orbits gives the common
quartic-block levels: this is Radeschi's composed foliation
\cite[Theorems~A--B and Section~3]{Radeschi14}. The one-leaf target
gives the OT--FKM foliation. In the splitting model the target
$\mathrm{SO}(m)$-action fixes $p_0$; its leaves retain $u-v$, while
the quartic subalgebra retains only $(u-v)^2$.
The balanced $(3,7,7)$ model in
\cref{sec:weighted-Clifford-quartic} uses $\mathcal C_{9,1}$ and
target $\mathrm{SO}(3)\times\mathrm{SO}(7)$-orbits, with profile
$E_+=E_-$. A profile equation selects a union of leaves; the
differential table requires no connectedness assumption.
\end{remark}

For three quartic blocks, the weights in the table are the target
block dimensions and the remaining weight $\ell-m$. The following
criterion identifies every minimal linear combination with a regular
zero meeting the positive part of this table.

\begin{proposition}[Minimal levels of orthogonal quartic generators]
\label{prop:orthogonal-quartics}
Suppose that nonzero homogeneous quartic polynomials $E_0,E_1,E_2$ satisfy
\begin{equation}\label{eq:orthogonal-quartic-data}
 E_0+E_1+E_2=R^2,\qquad
 \Gamma(E_i,E_j)=16RE_i\delta_{ij},\qquad
 \Delta E_i=8r_iR
\end{equation}
for constants $r_i>0$. They generate a Laplacian algebra together
with $R$, of generator degrees $(2,4,4)$ after eliminating one $E_i$.
Let $F=\sum_i\alpha_iE_i$ have a regular zero point with all $E_i>0$,
and put $\sigma=\sum_i r_i\alpha_i$.
Then $F=0$ is minimal if and only if its nonzero coefficients have
exactly two values $\alpha>0>\beta$, satisfying
\begin{equation}\label{eq:quartic-level-condition}
 \alpha(r_+-1)+\beta(r_--1)=0,\qquad
 r_+=\sum_{\alpha_i=\alpha}r_i,\quad r_-=\sum_{\alpha_i=\beta}r_i.
\end{equation}
If $F=0$ is minimal and none of the coefficients vanishes,
$\R[R,F]$ is already Laplacian.
\end{proposition}

\begin{proof}
The norm identities give $E_i\geq0$. Where they are positive,
their gradient matrix is positive definite, so their image has
interior. Set $A_j=\sum_i\alpha_i^jE_i$. The closure table gives
\[
 \Gamma(F,F)=16RA_2,\qquad
 \cL(F)=128R^2(\sigma A_2-A_3)-64FA_2.
\]
Thus minimality on the regular zero plane is equivalent to
$\sigma A_2-A_3=\kappa F$ for a constant $\kappa$. The image contains
an open piece of that plane, so comparison of its linear coefficients
is valid. For each nonzero $\alpha_i$ this reads
$\alpha_i(\sigma-\alpha_i)=\kappa$. There are at most two such
values, and a zero point with positive $E_i$ requires both signs.
Their sum is $\sigma$, which is precisely
\eqref{eq:quartic-level-condition}. These conditions also make
$\cL(F)$ divisible by $F$ globally and prove sufficiency.
In the minimal case, if all coefficients are nonzero, then
$A_2=(\alpha+\beta)F-\alpha\beta R^2$; the displayed norm and
Laplacian identities prove closure of $\R[R,F]$.
\end{proof}

\begin{remark}[More quartic blocks]\label{rem:quartic-many-blocks}
For any number of nonzero orthogonal quartic blocks summing to $R^2$,
the same coefficient comparison permits only two nonzero coefficient
values, provided the regular zero meets the positive-block region.
Group equal positive, negative, and zero coefficients and add their
weights. This reduces the minimal-level equation to
\eqref{eq:quartic-level-condition}, with two generators when no block
is unused. Different representations can still yield different cones.
\end{remark}

\begin{remark}[Compatible Cartan--M\"unzner equations]
\label{rem:compatible-quartic-CM}
Setting $\Phi_i=R^2-2E_i$ in \eqref{eq:orthogonal-quartic-data} gives
\[
 \sum_i\Phi_i=R^2,\qquad |\nabla\Phi_i|^2=16R^3,\qquad
 \Delta\Phi_i=4(N+2-4r_i)R.
\]
For distinct $i,j,k$ it also gives
$\Gamma(\Phi_i,\Phi_j)=-16R\Phi_k$.
Conversely, the sum and the individual gradient identities imply
these cross pairings; together with the displayed Laplacians they
recover \eqref{eq:orthogonal-quartic-data}.
These are compatible Cartan--M\"unzner-type equations, allowing
degenerate quartic representations: for example,
$E_i=(|x|^2-|y|^2)^2$ with equal-dimensional blocks has $r_i=1$.
\end{remark}

For the quartic algebra $\cA=\R[R,E_0,E_1,E_2]$ itself,
$\cA_2=\R R$. In the Clifford construction above, it lies in a
larger quadratic algebra. The criterion below determines its
realizable weights.

\begin{lemma}[Weights and image of the quartic blocks]
\label{lem:quartic-block-weights-image}
For the nonzero blocks in \eqref{eq:orthogonal-quartic-data}, the
weights are positive integers and $N=2(r_0+r_1+r_2-1)$. Their
spherical image is exactly
\begin{equation}\label{eq:quartic-simplex-image}
 (E_0,E_1,E_2)(\Sph^{N-1})
 =\{(u_0,u_1,u_2)\in[0,1]^3:u_0+u_1+u_2=1\}.
\end{equation}
\end{lemma}

\begin{proof}
Summing the Laplacian identities in \eqref{eq:orthogonal-quartic-data}
gives $N=2(r_0+r_1+r_2-1)$. The norm identities imply $E_i\geq0$.
On the unit sphere,
$u_i=E_i|_{\Sph^{N-1}}$ satisfies
\[
 |\nabla_{\Sph}u_i|^2=16u_i(1-u_i),\qquad
 \Delta_{\Sph}u_i=8r_i-4(N+2)u_i.
\]
A constant $u_i$ would be zero or one, contradicting the assumption
that all three blocks are nonzero.  Its minimum and maximum are
therefore zero and one.  Integration of the Laplacian identity
gives $r_i>0$.  At a nonzero zero of $E_i$, write
$H_i=\Hess E_i$.  Comparing quadratic terms in
$|\nabla E_i|^2=16RE_i$ gives
\[
 H_i^2=8RH_i,\qquad \operatorname{tr}H_i=8r_iR.
\]
Thus $H_i=8RP_i$ for an orthogonal projection $P_i$ of rank $r_i$;
in particular, each weight is an integer.

Nonnegativity and the sum of the blocks place the spherical image
in the displayed simplex. Its intersection with the simplex interior
is nonempty, since $E_0E_1E_2$ is a nonzero polynomial. There the
spherical gradient matrix
$16(\operatorname{diag}(u)-uu^{\mathsf T})$ has rank two, so the
intersection is relatively open by the submersion theorem. It is
also relatively closed, since the spherical image is compact.
Connectedness of the simplex interior, followed by taking its closure,
proves \eqref{eq:quartic-simplex-image}.
\end{proof}

The Clifford-module restriction gives an exact numerical range for
\eqref{eq:orthogonal-quartic-data}. We use the conventional function
\[
 (\delta(1),\ldots,\delta(8))=(1,2,4,4,8,8,8,8),
 \qquad \delta(k+8)=16\delta(k).
\]
Here $\delta(k)$ is the real irreducible module dimension for the
negative Clifford algebra on $k-1$ generators. A symmetric Clifford
system with $p\ge2$ matrices on $\R^{2\ell}$ exists exactly when
$\delta(p-1)\mid\ell$ \cite[Section~3.5]{FKM81}.
See also \cite[Chapter~0, pp.~3--4]{Shapiro2000} for the equivalent
Hurwitz--Radon restriction.

\begin{proposition}[Realizable quartic weights]
\label{prop:quartic-weight-realizability}
Let $r_0,r_1,r_2$ be positive integers, write their rearrangement as
$a\le b\le c$, and set $\ell=a+b+c-1$.
Nonzero quartics satisfying \eqref{eq:orthogonal-quartic-data} on
$\R^{2\ell}$ exist if and only if \(\delta(a+b-1)\mid\ell\).
Whenever this condition holds, a realization is obtained from one
Clifford system by two blocks of squared quadratic coordinates and
their complement in $R^2$.
\end{proposition}

\begin{proof}[Construction and parameter restrictions]
\leavevmode\par\nobreak
Assume the divisibility condition. Take $a+b$ symmetric Clifford
matrices, with quadratic coordinates $p_1,\ldots,p_{a+b}$.
The three polynomials
\[
 \sum_{\alpha=1}^{a}p_\alpha^2,\qquad
 \sum_{\alpha=a+1}^{a+b}p_\alpha^2,\qquad
 R^2-\sum_{\alpha=1}^{a+b}p_\alpha^2
\]
have weights $(a,b,c)$ by \cref{lem:Clifford-block-quartics}.
Their positive Laplacian coefficients ensure that none is zero.

For any realization, summing the Laplacians forces
$N=2(r_0+r_1+r_2-1)$. Positivity and integrality of the weights follow from
\cref{lem:quartic-block-weights-image}.
The converse uses the completed classification with four principal
curvatures and compatibility of the three blocks. See
\cref{app:quartic-weight-realizability} and its weight-one,
Clifford-joining, and exceptional-triple lemmas
\cref{lem:quartic-weight-one,lem:quartic-Clifford-joining,lem:quartic-exclude-155}.
\end{proof}

This is a classification of realizable weights, not of simultaneous
orthogonal equivalence classes of triples. For example, $(3,3,5)$
is excluded since $\delta(5)=8\nmid10$, although the individual
Cartan--M\"unzner multiplicity pairs $(2,7),(2,7),(4,5)$ all occur.
The criterion is symmetric in the weights; the subcalibration
range in \cref{thm:quartic-table-existence} instead depends on
which two blocks are chosen for the defining polynomial.

\begin{corollary}[The weights with active sum nine]
\label{cor:quartic-nine-weights}
Let $r_0,r_1,r_2$ be positive integers with $r_1,r_2\geq2$ and
$r_1+r_2=9$. A quartic realization exists if and only if $r_0$ is
a positive multiple of eight. Its ambient dimension is $N=2r_0+16$.
\end{corollary}

\begin{proof}
Up to exchanging the active weights, their pairs are $(2,7)$,
$(3,6)$, and $(4,5)$. For $1\leq r_0\leq7$, none of the corresponding module
dimensions divides $\ell=r_0+8$; the seven checks are listed in
\cref{tab:quartic-nine-low-weights}.
If $r_0\geq8$, the two
smallest weights are the active pair. Hence $a+b-1=8$ and the same
criterion is $8=\delta(8)\mid r_0+8$, equivalently $8\mid r_0$.
\end{proof}


\subsection{Subcalibrations from orthogonal quartic data}
\label{sec:quartic-table-existence}

The orthogonal-quartic table determines the minimal level, singular
set, and subcalibration of every realization in the following sufficient
minimizing range. Exact thresholds for particular families also use
the instability arguments below.

Throughout this subsection, let $E_0,E_1,E_2$ be nonzero homogeneous
quartic polynomials on $\R^N$ satisfying
\eqref{eq:orthogonal-quartic-data}, with $R=|x|^2$.  Suppose
$r_1,r_2\geq2$, and put
\begin{gather}
 \begin{gathered}
 h=r_1+r_2,\qquad \rho=r_0+r_1+r_2,\\
 c=\frac{r_1-1}{h-2},\qquad a=1-2c,\qquad b=c(1-c)>0,
 \end{gathered}
 \label{eq:quartic-general-weights}\\
 T=\frac{E_1+E_2}{4},\qquad
 F=\frac{(1-c)E_1-cE_2}{4}.
 \label{eq:quartic-general-generators}
\end{gather}
By \cref{lem:quartic-block-weights-image}, $N=2\rho-2$.
The choice of $c$ is the minimal-level condition
\eqref{eq:quartic-level-condition}.  Direct substitution gives
\begin{equation}\label{bh:reduced-table}
 \begin{gathered}
 (\Gamma(I_i,I_j))_{I=(R,T,F)}=
 \begin{pmatrix}
 4R&8T&8F\\
 8T&4RT&4RF\\
 8F&4RF&4R(aF+bT)
 \end{pmatrix},\\
 (\Delta R,\Delta T,\Delta F)=(2N,2hR,2aR).
 \end{gathered}
\end{equation}
Thus the gradient table is independent of $r_0$, while the ambient
dimension records all three weights.  By
\cref{prop:MR-generator-test}, $\R[R,T,F]$ is a Laplacian algebra.
It is exactly $\cA_R(F)$: writing $S=|\nabla F|^2$, one obtains
\[
 S=4R(aF+bT),\qquad
 \Delta S=8(N+8)(aF+bT)+8(a^2+bh)R^2.
\]
Since $b>0$, the second identity recovers $T$ in $\cA_R(F)$.
The generators are algebraically independent.  Indeed, the nonzero
polynomial $E_0E_1E_2$ is positive somewhere; there the gradients of
the three blocks are independent, and the change to $(R,T,F)$ is
locally invertible.  Fixing $R=1$ therefore leaves spherical quotient
dimension two.

\begin{lemma}[Geometry determined by the quartic data]
\label{lem:quartic-table-geometry}
At each fixed $R\geq0$, the full image of $(T,F)$ is
described by
\begin{equation}\label{eq:quartic-table-image}
 0\leq T\leq R^2/4,\qquad -cT\leq F\leq(1-c)T.
\end{equation}
The critical set of $F$ and the geometric singular set of
$C=\{F=0\}$ are both $Z=\{T=0\}$, and
$\dim Z\leq N-h$.  The hypersurface $C$ is minimal outside $Z$,
both phases approach every point of $C$, and
$\partial\llbracket\{F<0\}\rrbracket$ has support $C$.
\end{lemma}

\begin{proof}
The simplex image in \cref{lem:quartic-block-weights-image} and
homogeneity prove \eqref{eq:quartic-table-image}. Moreover,
\[
 |\nabla F|^2=R\bigl((1-c)^2E_1+c^2E_2\bigr),
\]
so the critical set of $F$ is exactly $Z$.  The reduced table gives
\begin{equation}\label{bh:minimality}
 \cL(F)=-8F\{b(R^2+2T)+2aF\},
\end{equation}
which proves regular minimality.

Nonzero points of $Z$ exist because $E_0$ attains one on the unit
sphere. Fix such a point $p$ and write $R_0=R(p)$ and
$\Hess E_i(p)=8R_0P_i$ for $i=1,2$, using the projections obtained
in the proof of \cref{lem:quartic-block-weights-image}. The quadratic terms of
$\Gamma(E_1,E_2)=0$ give $P_1P_2+P_2P_1=0$.  Taking the trace
shows $P_1P_2=0$.  Hence
\[
 \Hess T(p)=2R_0(P_1+P_2),\qquad
 \Hess F(p)=2R_0\bigl((1-c)P_1-cP_2\bigr).
\]
The first Hessian has rank $h$.  Applying the implicit function
theorem to $h$ independent components of $\nabla T$ places $Z$
locally in a smooth submanifold of codimension $h$.  Consequently
$\dim Z\leq N-h$ and $\cH^{N-1}(Z)=0$.
The second Hessian has positive and negative indices $r_1,r_2$,
respectively, so both phases occur arbitrarily near every such $p$.
If the zero set were a smooth hypersurface at $p$, its leading
quadratic term would vanish on its tangent hyperplane and thus have
rank at most two.  Its actual rank is $h\geq4$, a contradiction.
Thus these points are geometrically singular.

Homogeneity gives the same two-phase conclusion at the origin,
and regularity gives it on $C\setminus Z$.  Polynomial finite
perimeter and the local BV Poincar\'e inequality then identify the
support of the boundary current with $C$.  The vertex is singular
as well: a cone smooth at its vertex would be a hyperplane, whereas
$C$ has nonzero singular points.
\end{proof}

\begin{theorem}[Polynomial subcalibrations from orthogonal quartics]
\label{thm:quartic-table-existence}
For the data \eqref{eq:orthogonal-quartic-data} and
\eqref{eq:quartic-general-weights}--\eqref{eq:quartic-general-generators},
suppose that either
\[
 h\geq9,
 \qquad\text{or}\qquad
 r_1=r_2=4\ \text{ and }\ r_0\geq5.
\]
There exists a homogeneous polynomial $G=AF$ of degree at most ten,
with $A>0$ on $T>0$, whose critical set is $Z=\{T=0\}$ and whose
normalized gradient satisfies
\[
 \operatorname{sgn}(G)\,\diver X_G>0\qquad\text{off }C=\{F=0\}.
\]
The current $\partial\llbracket\{F<0\}\rrbracket$ has full support
$C$ and is strictly area minimizing.  Its regular part is strictly
stable, with the explicit lower bound
\begin{equation}\label{eq:bihomogeneous-Hardy}
 Q_C(\varphi)\geq
 \left[\frac{(N-17)^2}{4}+8h-72\right]
 \int_Cr^{-2}\varphi^2,\qquad
 \varphi\in C_c^\infty(C_{\reg}\setminus\{0\}),\quad r=\sqrt R.
\end{equation}
\end{theorem}

\subsubsection{Selecting the multiplier}

At $p\in Z\setminus\{0\}$, \eqref{bh:minimality} gives
$\lambda_F(p)=-8bR(p)^2<0$. Hence
\cref{lem:critical-multiplier-defect} forces every admissible smooth
multiplier to vanish on $Z$.
For $A\in\R[R,T]$, scaling a nonzero point of $Z$ gives
$A(R,0)=0$ for all $R>0$, so $T$ divides $A$ as a polynomial in
$R,T$.  The lowest two nonzero homogeneous choices are therefore
$T$ and $RT$, up to a positive constant.

These two candidates suffice in the stated ranges:
\begin{equation}\label{bh:potentials}
 G=\begin{cases}
 TF,&h\geq11,\ \text{or }h=10\ \text{and }r_1,r_2\geq3,\\
 RTF,&h=9,\ \text{or }\{r_1,r_2\}=\{2,8\},\\
 RTF,&r_1=r_2=4,\ r_0\geq5.
 \end{cases}
\end{equation}
With $s=T/R^2$ and $z=F/T$, these are $R^4s^2z$ and $R^5s^2z$:
the radial factor changes the defect but preserves vanishing on $Z$.
We use \cref{thm:finite-reduction,thm:exact-multiplier} and, for
$h=9$, the realizability restriction \cref{cor:quartic-nine-weights}.

\subsubsection{The degree-eight potential and its range}

For $G_8=TF$, the reduced table gives
\begin{gather}
 \norm{\nabla G_8}^2=4RT(bT^2+aTF+3F^2),\qquad
 \Delta G_8=R\{2aT+2(h+4)F\},\label{bh:G8-basic}\\
 \begin{split}\label{bh:G8-defect}
 \cL(G_8)=8TF\big\{&bT^2[(h-8)R^2-4T]\\
                  &+aTF[(h-4)R^2-4T]
                    +3F^2[hR^2-4T]\big\}.
 \end{split}
\end{gather}
Introduce the dimensionless variables
\[
 s=T/R^2\in(0,1/4],\qquad z=F/T\in[-c,1-c].
\]
The braces in \eqref{bh:G8-defect}, divided by $R^2T^2$, are
\begin{equation}\label{bh:p8}
 p_8(z,s)=b(h-8-4s)+a(h-4-4s)z+3(h-4s)z^2.
\end{equation}
Since
\[
 b+az=(1-c)^2(z+c)+c^2(1-c-z)>0,
\]
we have $\partial_s p_8=-4(b+az+3z^2)<0$.  It is enough to prove
\begin{equation}\label{bh:q8}
 q(z):=p_8(z,1/4)=b(h-9)+a(h-5)z+3(h-1)z^2>0.
\end{equation}

First take one of the degree-eight cases in \eqref{bh:potentials},
so $h\geq10$. Interchanging $E_1,E_2$ preserves the problem, so we
may assume
$1/(h-2)\leq c\leq1/2$ and $a\geq0$.  For $z\geq0$,
\eqref{bh:q8} is positive.  On $[-c,0]$ its vertex is
$z_*=-a(h-5)/(6(h-1))$, which lies in this interval exactly when
\[
 c\geq c_*:=\frac{h-5}{8(h-2)}.
\]
If $c<c_*$, the minimum is attained at $-c$, where \(q(-c)=4c((h-1)c-1)>0\).
If $c\geq c_*$, the minimum is positive exactly when
\begin{equation}\label{bh:D8}
 D_8(c):=12(h-1)(h-9)c(1-c)-(h-5)^2(1-2c)^2>0.
\end{equation}
This expression increases on $0<c\leq1/2$.  For $h\geq13$,
\[
 D_8(c_*)=
 \frac{3(h-1)(h-5)(h^2-14h+21)}{4(h-2)^2}>0.
\]
For $h=12,11$, the bound $c\geq1/(h-2)>c_*$ gives, respectively,
\[
 D_8(1/10)=107/25>0,\qquad D_8(1/9)=52/27>0.
\]
For $h=10$ with $r_1,r_2\geq3$, one has $c\geq1/4>c_*$ and
$D_8(1/4)=14>0$. This proves positivity in every degree-eight
case of \eqref{bh:potentials}.

The same candidate does not give a uniform result for $h=10$.
For $c=1/8$, hence $a=3/4$ and $b=7/64$, one has
\begin{equation}\label{eq:quartic-degree-eight-obstruction}
 p_8(-1/16,1/4)=-5/256<0.
\end{equation}
By \cref{lem:quartic-table-geometry}, this point is attained in every
realization with $(r_1,r_2)=(2,8)$. Thus $TF$ has the wrong divergence
sign at points of the complement. For $h=9$, one has
$q(z)=4az+24z^2$ with $a\ne0$, so $q$ is negative for small $z$ of
the opposite sign. For $r_1=r_2=4$, one instead has
$q(z)=-b+21z^2$, negative near zero. The exact image assertion again
realizes these negative values with $F\ne0$. Hence, within the
theorem's range, the degree-eight choices in \eqref{bh:potentials}
are precisely the cases in which $TF$ has the required strict sign.

\subsubsection{The degree-ten potential in the remaining ranges}

For $G_{10}=RTF$, the same table gives
\begin{equation}\label{bh:G10}
 \cL(G_{10})=8R^5T^3F\,p_{10}(z,s),
\end{equation}
where
\begin{align}
 p_{10}(z,s)={}&b[h-8+(2h-23)s]
       +a[h-4+(2h-14)s]z\nonumber\\
       &+[3h+(15h-33)s+(18h+27)s^2]z^2\nonumber\\
       &+2(r_0-1)s\bigl[b+az+3(1+3s)z^2\bigr].
       \label{bh:p10}
\end{align}
Because $r_0\geq1$ and $b+az>0$ on the image interval, the last
line is nonnegative.  Thus increasing the unused block weight
$r_0$ can only improve this certificate.

For $h=9$, \cref{cor:quartic-nine-weights} gives $r_0\geq8$.
At $r_0=8$, the full polynomial is
\begin{equation}\label{eq:quartic-nine-certificate}
 p_{10}(z,s)=b(1+9s)+a(5+18s)z+
 (27+144s+315s^2)z^2.
\end{equation}
Since $c\in[1/7,6/7]$, one has $a^2\leq25/49$ and
$4b\geq24/49$. The negative discriminant is therefore bounded below by
\begin{align}
 &4b(1+9s)(27+144s+315s^2)-a^2(5+18s)^2\nonumber\\
 &\qquad\geq
 \frac{23+4788s+30564s^2+68040s^3}{49}>0.
 \label{eq:quartic-nine-discriminant}
\end{align}
The leading coefficient is positive, so \eqref{eq:quartic-nine-certificate}
is positive for all real $z$ and $0<s\leq1/4$. Increasing $r_0$
preserves positivity on the image interval.

When $h=10$, the first three terms are
\[
 b(2-3s)+a(6+6s)z+(30+117s+207s^2)z^2.
\]
Here $c\in[1/8,7/8]$, so $a^2\leq9/16$ and
$b=(1-a^2)/4$.  The negative of the discriminant is bounded below by
\begin{align*}
 &4b(2-3s)(30+117s+207s^2)-a^2(6+6s)^2\\
 &\hspace{10mm}\geq
 6+\frac{45}{2}s+\frac{117}{16}s^2-\frac{4347}{16}s^3
 \geq\frac{1797}{1024}>0.
\end{align*}
The final inequality follows from $0<s\leq1/4$ by discarding the
positive linear and quadratic terms. The leading coefficient is
positive, so this polynomial is positive for all real $z$ and all
allowed $s$.

Finally, for $r_1=r_2=4$ one has $h=8$, $a=0$, and $b=1/4$.
The full polynomial becomes
\begin{equation}\label{eq:quartic-balanced-eight-certificate}
 p_{10}(z,s)=\frac{2r_0-9}{4}s+
 \bigl[24+(81+6r_0)s+(153+18r_0)s^2\bigr]z^2.
\end{equation}
This is positive for every $s>0$ if $r_0\geq5$.
Together these calculations prove the required sign in every
case of \eqref{bh:potentials}.

\subsubsection{Global minimality and the Hardy bound}

\begin{proof}[Completion of the proof of \cref{thm:quartic-table-existence}]
\leavevmode\par\nobreak
Both multipliers are positive on $T>0$. Euler's identity excludes
critical points where $F\ne0$, and $|\nabla F|^2=4bRT$ excludes them
on $F=0$, $T>0$. On $Z$, both potentials and their gradients vanish.
Thus they have critical set $Z$ and the same phases as $F$.
The full-support and dimension conclusions of
\cref{lem:quartic-table-geometry} verify the hypotheses of
\cref{cor:gradient-subcalibration,thm:removable-subcalibration}.
Their comparison, followed by \cref{cor:polynomial-strict-minimality},
gives strict integral-current mass minimization and the gap
$\Theta\varepsilon^{N-1}$.

For stability, \eqref{eq:boundary-potential-algebraic} gives
\[
 \beta_{TF}=2(h-8)\frac{R^2}{T}-8,\qquad
 \beta_{RTF}=2(h-8)\frac{R^2}{T}+4\rho-50.
\]
Since $R^2/T\geq4$ and $h\geq8$, the homogeneous Hardy estimate
\cref{thm:ground-state-gap}, with cone dimension $N-1=2\rho-3$,
gives the same lower bound for both potential degrees:
\begin{align*}
 \left(8-\frac{N-1}{2}\right)^2+8(h-8)-8
 &=\frac{(N-17)^2}{4}+8h-72,\\
 \left(10-\frac{N-1}{2}\right)^2+8(h-8)+4\rho-50
 &=\frac{(N-17)^2}{4}+8h-72.
\end{align*}
For $h\geq10$ this coefficient is positive. For $h=9$,
\cref{cor:quartic-nine-weights} gives $N=2r_0+16\geq32$, so the
coefficient $(N-17)^2/4$ is positive as well. In the remaining
case, $h=8$ and $r_0\geq5$ imply $N\geq24$, so it is at least
$17/4$.  This proves \eqref{eq:bihomogeneous-Hardy}.
\end{proof}

For the splitting specialization $r_0=1$, $h=\ell$, one has
$N=2\ell$, and the Hardy coefficient becomes
$\ell^2-9\ell+1/4$. For $r_0>1$ and unequal active weights,
\cref{prop:quartic-weight-realizability} determines realizability,
and the inequalities above give a sufficient minimizing range.

\subsubsection{The balanced specialization}

\begin{corollary}[The degree-eight potential for balanced quartics]
\label{prop:balanced-quartic-multiplier}
Suppose that the three nonzero quartics in
\eqref{eq:orthogonal-quartic-data} have weights $(r_0,\nu,\nu)$,
with $\nu\geq5$. Set
\[
 F=\frac{E_+-E_-}{4},\qquad W=\frac{E_++E_-}{4},\qquad G=WF.
\]
The conclusions of \cref{thm:quartic-table-existence} hold for $G$,
with $\Sigma=\{W=0\}$. In this normalization, $\Delta F=0$,
$|F|\leq W\leq R^2/4$, and
\begin{equation}\label{eq:balanced-quartic-smallest-algebra}
 \cA_R(F)=\R[R,W,F]
\end{equation}
is a polynomial algebra with generator degrees $(2,4,4)$.
Moreover,
\begin{align}
 |\nabla G|^2&=4RW(W^2+3F^2),\qquad
 \Delta G=4(\nu+2)RF,\label{eq:balanced-quartic-gradient}\\
 \cL(G)&=16FW\left\{
 W^2\big[(\nu-4)R^2-2W\big]
 +3F^2\big[\nu R^2-2W\big]\right\}.
 \label{eq:balanced-quartic-certificate}
\end{align}
\end{corollary}

\begin{proof}
Take $T=W$, $F_{\mathrm{gen}}=F/2$, $h=2\nu\geq10$, and $c=1/2$
in \cref{thm:quartic-table-existence}. Its table becomes
\begin{equation}\label{eq:balanced-quartic-table}
 \Delta(R,W,F)=(2N,4\nu R,0),\qquad
 (\Gamma(I_i,I_j))_{I=(R,W,F)}=
 \begin{pmatrix}
 4R&8W&8F\\
 8W&4RW&4RF\\
 8F&4RF&4RW
 \end{pmatrix}.
\end{equation}
Substituting this normalization in
\eqref{bh:G8-basic}--\eqref{bh:G8-defect} gives
\eqref{eq:balanced-quartic-gradient}--\eqref{eq:balanced-quartic-certificate}.
\end{proof}

The following table organizes the applications of the common
quartic theorem. The splitting, Riedler, and weight-two families
have exact thresholds, obtained by combining this theorem with
instability arguments. Each construction also identifies its
singular set, including the real rank-one strata in the weight-two
families.

\begin{table}[htbp]
\centering\small
\setlength{\tabcolsep}{5pt}
\renewcommand{\arraystretch}{1.15}
\caption{Euclidean quartic realizations of
\cref{thm:quartic-table-existence}. Only parameters admitting the
specified Clifford realization are included. The intrinsic
Simons-cone result in \cref{thm:Hopf-Simons-section} uses a
separate differential table.}
\label{tab:quartic-weight-realizations}
\begin{tabular}{@{}lll@{}}
\toprule
Family & $(r_0,r_1,r_2)$ & Minimizing parameters used here\\
\midrule
OT--FKM splitting & $(1,m,\ell-m)$ & $\ell\ge10$, $2\le m\le\ell-2$\\
Riedler pullback & $(N/2-2k+1,k,k)$ & $N\ge32$, $k\ge4$\\
Hopf pairing & $(1,8,8)$ & Ambient dimension $32$\\
Balanced Clifford example & $(3,7,7)$ & Ambient dimension $32$\\
Additional block example & $(2,5,10)$ & Ambient dimension $32$\\
Weight-two families & $(2,m,4k-m-1)$ & $m=2,3$, $k\ge3$\\
Unequal Lawson pullback & $(22,5,6)$ & Ambient dimension $64$\\
\bottomrule
\end{tabular}
\end{table}

\subsection{The splitting, Hopf and Riedler families}
\label{sec:quartic-family-completions}

We apply the common theorem to the splitting construction, the Hopf
pairing and Riedler's pullbacks. The splitting and Riedler families
admit complete variational classifications. The balanced Hopf section
also illustrates intrinsic comparison inside a Simons cone, and the
degree-nine Riedler potential supplies an additional Hardy improvement.

\subsubsection{The OT--FKM splitting criterion}

For the data \eqref{bh:quartic-data}--\eqref{bh:centered}, the exact
minimizing threshold is as follows.

\begin{theorem}\label{bh:main}
The hypersurface $C=\{F=0\}\subset\R^{2\ell}$ is minimal on its
regular part, and
\[
 \Sing C=(\R^\ell\times\{0\})\cup(\{0\}\times\R^\ell).
\]
The oriented boundary $\partial\llbracket\{F<0\}\rrbracket$ is
area minimizing if and only if $\ell\geq10$. In all these minimizing
cases it is strictly area minimizing in the sense of
\cref{def:four-variational-notions}. The polynomial potentials and
strict-stability bound of \cref{thm:quartic-table-existence} apply,
with degrees ten for $\ell=10$ and eight for $\ell\geq12$.
For the remaining
admissible dimensions $\ell=4,6,8$, the regular part of $C$ is unstable.
\end{theorem}

\paragraph{Closure and singularities.}

For $x\ne0$, the vectors \(x,A_1x,\ldots,A_{m-1}x\)
are orthogonal and have length $|x|$. Consequently $0\leq Q\leq T$.
Direct differentiation gives, in the variables $I=(u,v,Q)$,
\begin{equation}\label{bh:original-table}
 (\Gamma(I_i,I_j))=
 \begin{pmatrix}4u&0&4Q\\0&4v&4Q\\4Q&4Q&4RQ\end{pmatrix},
 \qquad
 \Delta I=(2\ell,2\ell,2mR).
\end{equation}
For example,
$\nabla_yQ=2\sum_{i=0}^{m-1}\ip{A_i x}{y}A_i x$, with $A_0=\Id$,
so $\norm{\nabla_yQ}^2=4uQ$; differentiating in $x$ gives $4vQ$.
By \cref{prop:MR-generator-test},
\begin{equation}\label{bh:algebra}
 \cA=\R[u,v,Q]=\R[R,u-v,F]
\end{equation}
is a graded Laplacian algebra, with generator degrees $(2,2,4)$.
The determinant of its gradient matrix is $64RQ(T-Q)$, which is
positive at points with $u,v>0$ and $0<Q<T$. These points exist by
choosing the components of $y$ both in and perpendicular to the
span above. Thus the three generators are algebraically independent;
the spherical quotient dimension is two. The separately harmonic
normal form of \cref{prop:two-block-normal} is
$P=Q-(m/\ell)uv$; the minimal polynomial $F=Q-cuv$ differs from
it by a radial term in the two blocks.

With
$p_0=u-v$ and $p_{i+1}=2\langle A_i x,y\rangle$ for
$A_0=\Id$, the matrices in \cref{rem:split-composed-foliation}
form a symmetric Clifford system, and
\[
 F=\frac14\left(\sum_{i=1}^{m}p_i^2
                    +cp_0^2-cR^2\right).
\]
Its blocks $E_0=p_0^2$, $E_1=4Q$, and $E_2=4(T-Q)$
have weights $(1,m,\ell-m)$. The substitutions in
\eqref{bh:centered} therefore give \eqref{bh:reduced-table} with
$h=\ell$ and $N=2\ell$.
The calculation after \eqref{bh:reduced-table} recovers
$\cA_R(F)=\R[R,T,F]$, of degrees $(2,4,4)$. Independence follows
since $(u,v,Q)\mapsto(u+v,uv,Q-cuv)$ has Jacobian $u-v$.
Both algebras have spherical quotient dimension two, but
\eqref{bh:algebra} also retains the sign of $u-v$.

The geometry follows from \cref{lem:quartic-table-geometry}. For
$\ell\ge10$, evenness of $\ell$ and
\cref{thm:quartic-table-existence} give the variational conclusions.
The remaining dimensions are excluded by the following local
model and angular calculation.

Fix $x_0\ne0$, and write
$W_{x_0}=\operatorname{span}\{x_0,A_1x_0,\ldots,A_{m-1}x_0\}$.
The first nonzero term of $F$ at $(x_0,0)$ is
\begin{equation}\label{bh:transverse}
 \norm{x_0}^2\big((1-c)\norm{y_\parallel}^2
                         -c\norm{y_\perp}^2\big),
 \quad y_\parallel\in W_{x_0},\quad y_\perp\in W_{x_0}^\perp.
\end{equation}
The tangent cone is the product of $\R^\ell$ with this nonlinear
Lawson cone. Its transverse link is
\[
 \Sph^{m-1}(\sqrt c)\times
 \Sph^{\ell-m-1}(\sqrt{1-c})\subset\Sph^{\ell-1}.
\]
Both multiplicities are positive by \eqref{bh:range}. Smooth local
orthonormal frames for $W_x$ and $W_x^\perp$ identify a neighborhood
of the stratum with a smoothly varying family of these quadratic cones.
The same description holds on the other coordinate subspace. 

\begin{remark}[Clifford coordinates of the OT--FKM splitting]
\label{rem:split-composed-foliation}
Put $A_0=\Id$. The symmetric Clifford system
\[
 P_0=\begin{pmatrix}\Id&0\\0&-\Id\end{pmatrix},\qquad
 P_{i+1}=\begin{pmatrix}0&A_i^T\\A_i&0\end{pmatrix}
 \quad(0\leq i\leq m-1)
\]
has coordinates $p_0=u-v$ and $p_{i+1}=2\langle A_i x,y\rangle$.
The common geometric interpretation is given in
\cref{sec:Clifford-composition-background}.
\end{remark}


\subsubsection{The unstable OT--FKM splitting dimensions}

For $\ell=4,6$, the transverse Lawson cone in \eqref{bh:transverse}
is unstable: its link has dimension $\ell-2$ and constant squared
curvature $\ell-2$, giving the Hardy margin
\begin{equation}\label{bh:Lawson-margin}
 \frac{(\ell-3)^2}{4}-(\ell-2)<0.
\end{equation}
\Cref{lem:tangent-instability-transfer} transfers a negative
variation from this transverse cone to $C_{\rm reg}$, using the
local model \eqref{bh:transverse}.

For $\ell=8$, the transverse Lawson cone can be area minimizing;
instability follows instead from the angular calculation below.

\begin{lemma}\label{bh:curvature}
On the regular part of the quartic cone,
\begin{equation}\label{bh:AC}
 \norm{A_C}^2=(\ell-2)\frac{R}{T}+\frac{2(\ell-1)}R.
\end{equation}
On $\Lambda=C\cap\Sph^{2\ell-1}$, put $t=T|_\Lambda$. Then
\begin{gather}\label{bh:link-table}
 \norm{A_\Lambda}^2=\frac{\ell-2}{t}+2(\ell-1),\qquad
 \norm{\nabla_\Lambda t}^2=4t(1-4t),\\
 \Delta_\Lambda t=2(\ell-1)-8\ell t.\nonumber
\end{gather}
\end{lemma}
\begin{proof}
Put $W=\norm{\nabla F}^2$ and $\nu=\nabla F/\sqrt W$.
The Euclidean Bochner identity gives
$\norm{\nabla^2F}^2=\Delta W/2-\Gamma(F,\Delta F)$.
On $F=0$, \eqref{bh:reduced-table} yields
\begin{align*}
 W&=4bRT,\\
 \norm{\nabla^2F}^2&=4(a^2+b\ell)R^2+8b(\ell+4)T,\\
 \norm{\nabla\sqrt W}^2&=4[(a^2+b)R^2+5bT],\\
 \nabla^2F(\nu,\nu)&=\Delta F=2aR.
\end{align*}
Now use
\[
 \norm{A_C}^2=
 \frac{\norm{\nabla^2F}^2-2\norm{\nabla\sqrt W}^2
                     +(\nabla^2F(\nu,\nu))^2}{W}
\]
to obtain \eqref{bh:AC}. The normal Hessian of any polynomial $V$ can also be computed
from the closure table, using
\[
 \Hess V(\nu,\nu)=\frac{
 \Gamma(F,\Gamma(F,V))-\tfrac12\Gamma(V,W)}{W}.
\]
On $F=0$, this gives
\[
 \Gamma(T,F)=0,\qquad
 \Hess T(\nu,\nu)=2R-\frac{4T}{R},\qquad
 \Delta_C T=2(\ell-1)R+\frac{4T}{R}.
\]
Subtracting the radial contribution of the degree-four function $T$
on the $(2\ell-1)$-dimensional cone gives the last identity in
\eqref{bh:link-table}. Subtracting the squared radial derivative from
$\Gamma(T,T)=4RT$ gives the gradient identity.
\end{proof}

For $\ell=8$, the positive function $\phi=t^{-1}$ satisfies
\begin{equation}\label{bh:eigenfunction}
 (\Delta_\Lambda+\norm{A_\Lambda}^2)\phi=46\phi.
\end{equation}
To apply this identity, we cut off near the singular set of the link.
Its two singular strata have dimension seven and
transverse dimension seven in $\Lambda$. The local quadratic-cone
description after \eqref{bh:transverse} implies
$t\asymp\rho^2$ and transverse volume density comparable to
$\rho^6\,d\rho$, where $\rho$ is distance to the singular stratum.
Thus $\phi\asymp\rho^{-2}$, its gradient is $O(\rho^{-3})$, and
$\phi$, its gradient and the potential term have the required
finite integrals.

Choose cutoffs $\eta_\varepsilon$ vanishing for $\rho<\varepsilon$
and equal to one for $\rho>2\varepsilon$, with gradient bounded by
$O(\varepsilon^{-1})$. Then
\[
 \int_\Lambda\phi^2\norm{\nabla\eta_\varepsilon}^2
 =O(\varepsilon^{-4}\varepsilon^{-2}\varepsilon^7)
 =O(\varepsilon)\longrightarrow0.
\]
Multiplying \eqref{bh:eigenfunction} by
$\eta_\varepsilon^2\phi$ and integrating by parts on the regular
part proves that the angular Rayleigh quotients of
$\eta_\varepsilon\phi$ tend to $-46$. The cone has dimension fifteen,
so \eqref{eq:exact-Hardy-gap} gives compactly supported regular
variations whose scale-invariant quotients tend to
\begin{equation}\label{bh:negative-margin}
 \frac{(15-2)^2}{4}-46=-\frac{15}{4}<0.
\end{equation}
This proves instability. Together with \eqref{bh:Lawson-margin} and
the evenness of $\ell$, it completes Theorem~\ref{bh:main}.

\subsubsection{The Hopf pairing and its balanced section}
\label{sec:Hopf-pairing}

The Hopf pairing realizes weights $(1,8,8)$ and has two linear
singular strata. Its balanced section also gives an intrinsic
comparison inside a Simons cone.

Let \(h:\R^{16}\to\R^9\) be the quadratic extension of the standard
octonionic Hopf map \(\Sph^{15}\to\Sph^8\), normalized by
\begin{equation}\label{eq:Hopf-normalization}
 |h(x)|^2=|x|^4,\qquad
 Dh(x)Dh(x)^{\mathsf T}=4|x|^2\Id_9,\qquad \Delta h=0.
\end{equation}
For polarization of the spherical norm identity and the Hopf construction
from norm-preserving bilinear multiplication, see
\cite[Chapter~15]{Shapiro2000}.
The coordinates are \(h_i(x)=\langle A_ix,x\rangle\) for a nine-matrix Clifford
system on \(\R^{16}\); see \cite[Theorem~A and Section~1.2]{Radeschi14}.
On \(\R^{16}\oplus\R^{16}\), set
\begin{equation}\label{eq:Hopf-pairing-definitions}
 u=|x|^2,\quad v=|y|^2,\quad R=u+v,\quad V=u-v,\quad T=uv,
 \qquad F(x,y)=\langle h(x),h(y)\rangle.
\end{equation}
In particular, \(|F|\leq T\), and \(F\) has bidegree \((2,2)\).

\begin{proposition}[The pairing closure]\label{prop:Hopf-closure}
The algebra \(\cA=\R[u,v,F]=\R[R,V,F]\) is a graded Laplacian
algebra of spherical quotient dimension two. Its closure table, in
the ordered generators \((u,v,F)\), is
\begin{equation}\label{eq:Hopf-closure}
 (\Delta u,\Delta v,\Delta F)=(32,32,0),\qquad
 \bigl(\Gamma(I_i,I_j)\bigr)=
 \begin{pmatrix}
  4u&0&4F\\
  0&4v&4F\\
  4F&4F&4uv(u+v)
 \end{pmatrix}.
\end{equation}
The quartic zero set \(C=\{F=0\}\) is minimal on its regular part,
and its geometric singular set is
\begin{equation}\label{eq:Hopf-singular-set}
 Z=(\R^{16}\times\{0\})\cup(\{0\}\times\R^{16}).
\end{equation}
\end{proposition}

\begin{proof}
Equation \eqref{eq:Hopf-normalization} gives
\[
 |\nabla_xF|^2=4uv^2,\qquad
 |\nabla_yF|^2=4u^2v,\qquad \Delta F=0.
\]
Euler's identity in each block gives the remaining entries of
\eqref{eq:Hopf-closure}, and \cref{prop:MR-generator-test}
proves Laplacian closure. Surjectivity of the Hopf map gives
the full image
\[
 \{(u,v,f):u,v\geq0,\ |f|\leq uv\}.
\]
Its interior proves algebraic independence and spherical quotient
dimension two. Since $uv$ varies at fixed $R$ and $F=0$,
$\Gamma(F,F)=4Ruv\notin\R[R,F]$.

Since \(\Gamma(F,R)=8F\) and \(\Gamma(F,T)=4RF\), we obtain
\[
 \cL(F)=-\tfrac12\Gamma(F,4RT)=-8(R^2+2T)F.
\]
This proves regular minimality and identifies the critical set as
\(Z\). At \((0,y_0)\), \(y_0\ne0\), the transverse quadratic term
is \(\langle A(h(y_0))x,x\rangle\), where
\(A(a)=\sum_i a_iA_i\). Its eigenvalues are
\(\pm|y_0|^2\), each with multiplicity eight. Hence the tangent
cone is the product of a quadratic cone of signature \((8,8)\)
with \(\R^{16}\), and the point is geometrically singular.
The other axis is identical; the vertex is singular as well.
\end{proof}

\begin{remark}[A second realization of the pairing table]
The Clifford construction in \cref{sec:Clifford-composition-background}
also realizes this table. Put
\(Q=u^2+v^2-2F\). For a nine-matrix Clifford system
\(P_0,\ldots,P_8\) on \(\R^{32}\), the generators
\[
 \widehat R=|z|^2,\qquad
 \widehat V=\langle P_0z,z\rangle,\qquad
 \widehat Q=\sum_{i=0}^8\langle P_iz,z\rangle^2
\]
have the same gradient and Laplacian identities as \((R,V,Q)\).
Both images are
\(\{(R,V,Q):R\geq0,\ V^2\leq Q\leq R^2\}\).
The inverse relations
\(u=(R+V)/2\), \(v=(R-V)/2\), and \(F=(R^2+V^2-2Q)/4\)
therefore identify their graded differential algebras.
\end{remark}

\begin{theorem}[A polynomial subcalibration of the pairing cone]
\label{thm:Hopf-pairing}
The potential \(G=TF=|x|^2|y|^2F\) has degree eight and
subcalibrates \(C=\{F=0\}\). Its oriented boundary current
\(\partial\llbracket\{F<0\}\rrbracket\) is strictly area minimizing
in \(\R^{32}\). Moreover,
\begin{equation}\label{eq:Hopf-Hardy-bound}
 Q_C(\varphi)\geq\frac{449}{4}\int_C r^{-2}\varphi^2,
 \qquad \varphi\in C_c^\infty(C_{\reg}),\quad r=\sqrt R.
\end{equation}
\end{theorem}

\begin{proof}
The three blocks
\[
 E_0=(u-v)^2,\qquad E_1=2(T+F),\qquad E_2=2(T-F)
\]
have weights $(1,8,8)$. The centered generator $F/2$ gives
$h=16$, $r_0=1$, and $c=1/2$ in
\cref{thm:quartic-table-existence}. Its balanced specialization
\cref{prop:balanced-quartic-multiplier} gives $G=TF$ and
\begin{align}
 \Delta G&=40RF,\qquad
 |\nabla G|^2=4RT(T^2+3F^2),\label{eq:Hopf-gradient}\\
 \cL(G)&=32TF\bigl[T^2(2R^2-T)
                    +3F^2(4R^2-T)\bigr].\label{eq:Hopf-defect}
\end{align}
The Hardy constant in \eqref{eq:bihomogeneous-Hardy} is $449/4$
at $(N,h)=(32,16)$.
\end{proof}

\begin{remark}[Relation with the OT--FKM and Riedler constructions]
\label{rem:Hopf-FKM-Riedler}
The matrices \(B_i=\operatorname{diag}(A_i,-A_i)\) form a Clifford
system on \(\R^{32}\). Its OT--FKM polynomial
\cite[Section~1.3]{Radeschi14} is
\[
 \Psi=R^2-2\sum_{i=1}^9\langle B_i(x,y),(x,y)\rangle^2
     =4F-V^2.
\]
It has multiplicities \((8,7)\). Applying the centering formula
\eqref{eq:iso-F-z} to \(-\Psi\), with ordered multiplicities
\((7,8)\), shows that its minimal regular hypercone is
\(\{\Psi+R^2/15=0\}\). Nonvertex singularities distinguish $C$ from every cone over a
regular isoparametric hypersurface. They also distinguish it from
Riedler's family: its singular dimension is $N-2k$, so congruence
would require $N=32$, $2k=16$, but sixteen Clifford matrices
require ambient dimension at least $256$
\cite[Remark~5.2.6(1) and Table~4.1]{RiedlerThesis25}. Although
\(F=\langle\cdot,\cdot\rangle\circ(h,h)\), the product map
\(\Phi(x,y)=(h(x),h(y))\) satisfies
\[
 D\Phi D\Phi^{\mathsf T}=4\operatorname{diag}(u\Id_9,v\Id_9).
\]
It is not horizontally conformal on \(u\ne v\), so the pullback
criterion in \cref{prop:pullback-subcalibration} does not apply
to this map on \(\R^{32}\).
\end{remark}

\paragraph{The balanced section inside a Simons cone.}
\label{sec:Hopf-intrinsic-extension}

For the pairing polynomial above, the ambient space for comparison
is now itself a minimal cone. Let
\[
 \mathcal B=\{V=0\}=\{|x|=|y|\},\qquad
 D=\{V=F=0\},\qquad E_{\mathcal B}=\mathcal B\cap\{F<0\}.
\]
The Simons cone \(\mathcal B\) is a singular space of dimension
\(31\), smooth away from the origin. We use its induced metric
and the orientation determined by \(\nabla V\).
Intrinsic perimeter is variation in the smooth part, extended across
the vertex by the cutoff below. Write it as \(\Per_{\mathcal B}\), and set
\(D_a=\partial\llbracket E_{\mathcal B}\rrbracket
\mathbin{\llcorner}B_a(0)\).

\begin{theorem}[The pairing section of a Simons cone]
\label{thm:Hopf-Simons-section}
The cone \(D\) has dimension \(30\), is smooth away from the
origin, and is minimal in \(\R^{32}\). The intrinsic potential
\(G_{\mathcal B}=RF\) subcalibrates it in \(\mathcal B\).
For every \(a>0\), the integral current \(D_a\) minimizes mass
among integral \(30\)-currents supported in
\(\mathcal B\cap\overline B_a(0)\) with boundary \(\partial D_a\).
Thus \(D_a\) solves the Plateau problem in the singular ambient
space \(\mathcal B\), with its Euclidean induced area.
\end{theorem}

\begin{proof}
On \(\mathcal B\setminus\{0\}\), the gradients of \(V\) and
\(F\) are nonzero and orthogonal by \eqref{eq:Hopf-closure}.
Hence \(D\setminus\{0\}\) is smooth of codimension two.
The unit normal to \(\mathcal B\) is
\(n_{\mathcal B}=(x,-y)/\sqrt R\). The bidegree of \(F\) gives
\(\Hess F((x,-y),(x,-y))=-4F\). Since \(\mathcal B\) is minimal,
the induced operators satisfy
\begin{equation}\label{eq:Hopf-intrinsic-closure}
 \begin{gathered}
 \Delta_{\mathcal B}R=62,\qquad
 \Delta_{\mathcal B}F=4F/R,\\
 \Gamma_{\mathcal B}(R,R)=4R,\qquad
 \Gamma_{\mathcal B}(R,F)=8F,\qquad
 \Gamma_{\mathcal B}(F,F)=R^3.
 \end{gathered}
\end{equation}
Here \(\nabla F\) is tangent to \(\mathcal B\), and \(T=R^2/4\).
The same chain rule as in \cref{thm:finite-reduction} now yields
\begin{equation}\label{eq:Hopf-intrinsic-defect}
 \begin{gathered}
 \Delta_{\mathcal B}(RF)=82F,\qquad
 |\nabla_{\mathcal B}(RF)|^2=R^5+20RF^2,\\
 \cL_{\mathcal B}(RF)=16RF(2R^4+85F^2),
 \end{gathered}
\end{equation}
where \(\cL_{\mathcal B}\) is \eqref{eq:L-definition} for the
induced metric. The normalized gradient
\[
 X=\frac{\nabla_{\mathcal B}(RF)}{|\nabla_{\mathcal B}(RF)|}
\]
is smooth on \(\mathcal B\setminus\{0\}\), equals the exterior
normal along \(D\setminus\{0\}\), and its divergence has the
strict sign of \(F\) off \(D\).

We apply the proof of \cref{thm:removable-subcalibration} in the
smooth part of \(\mathcal B\). Only the vertex requires a cutoff.
For a radial cutoff \(\eta_\eps\), zero on \(B_\eps(0)\) and
one outside \(B_{2\eps}(0)\),
\[
 \int_{\mathcal B\cap B_{2\eps}(0)}
       |\nabla_{\mathcal B}\eta_\eps|\,d\cH^{31}
 =O(\eps^{30}).
\]
Hence extending a bounded characteristic function across the vertex
adds no perimeter mass. The estimate $|\diver_{\mathcal B}X|\le c/r$
from homogeneity makes the divergence error vanish as well. For every relatively compact open set \(U\subset\mathcal B\)
and every locally finite-perimeter set \(H\subset\mathcal B\) with
\(H\triangle E_{\mathcal B}\Subset U\),
\begin{equation}\label{eq:Hopf-intrinsic-comparison}
 \Per_{\mathcal B}(H;U)-\Per_{\mathcal B}(E_{\mathcal B};U)
 \geq\int_{H\triangle E_{\mathcal B}}
       |\diver_{\mathcal B}X|\,d\cH^{31}\geq0.
\end{equation}
The set \(E_{\mathcal B}\) has locally finite perimeter because
its boundary is a cone with a smooth compact link.

For the asserted integral-current comparison, cone a competitor
minus \(D_a\) to the origin inside \(\mathcal B\cap\overline B_a(0)\).
This gives an integral \(31\)-dimensional filling. On the oriented smooth part of \(\mathcal B\), write this filling as
\(\llbracket k\rrbracket_{\mathcal B}\), with integer-valued
locally BV multiplicity, and set
\(v=\one_{E_{\mathcal B}}+k\) and \(H=\{v\geq1\}\).
The integer-level decomposition used in
\cref{thm:strict-minimality-flux} gives
\(\lvert D\one_H\rvert\leq|Dv|\) on the smooth part.
Writing \(T\) for the competitor, for \(b>a\) we obtain
\[
 0\leq \Per_{\mathcal B}(H;B_b)-
          \Per_{\mathcal B}(E_{\mathcal B};B_b)
 \leq \mathbf M(T)-\mathbf M(D_a).
\]
The inequalities use \eqref{eq:Hopf-intrinsic-comparison} and
cancellation of the common exterior mass; $D$ has no mass on a sphere.
The representation and integer-level decomposition hold in smooth
charts. The vertex cutoff extends the characteristic function,
and no positive-dimensional integral current is supported at the
vertex. Thus the comparison includes competitors meeting it.

Finally, the divergence in \eqref{eq:Hopf-intrinsic-defect} vanishes
on \(D\), so \(D\) is minimal in \(\mathcal B\).
For its unit normal \(\nu\) within \(\mathcal B\), the two block
lengths satisfy \(|\nu_x|=|\nu_y|\), by the partial gradient
identities above and \(u=v\). Hence
\(\mathrm{II}_{\mathcal B}(\nu,\nu)=0\). Since
\(\operatorname{tr}_{T\mathcal B}\mathrm{II}_{\mathcal B}=0\),
its trace on \(TD\) is also zero. The second fundamental form
composition formula proves Euclidean minimality of \(D\setminus\{0\}\).
The smooth compact link gives
\(\int_{D\cap B_{2\eps}}|\nabla_D\eta_\eps|\,d\cH^{30}
=O(\eps^{29})\to0\), so the Euclidean first variation also vanishes
across the vertex.
\end{proof}

The Euclidean area-minimizing cones of codimension two in
\cite[Theorem~B]{CuiFKM25} are also area-minimizing hypersurfaces within
the Simons cones containing them. Indeed, let \(K\subset\mathcal S\)
be such a cone and put \(K_a=K\mathbin{\llcorner}B_a(0)\).
Every integral current \(T\) supported in
\(\mathcal S\cap\overline B_a(0)\) with
\(\partial T=\partial K_a\) is also an admissible Euclidean
competitor, so \(\mathbf M(K_a)\leq\mathbf M(T)\).
The induced area on \(\mathcal S\) is the Euclidean area.




\subsubsection{Riedler's Clifford--Simons quartics}

The construction and area-minimizing range
\(N\geq32\), \(k\geq4\) are due to Riedler
\cite[Section~5.2 and Theorem~5.2.5]{RiedlerThesis25}.
The common quartic theorem gives polynomial potentials; the
degree-nine potential below improves their Hardy bound. The
sixteen-dimensional endpoint is a splitting cone.

Let \(A_1,\ldots,A_{2k}\) be a symmetric Clifford system on
\(\R^N\), and write
\(P_i(x)=\langle A_i x,x\rangle\) and
\(Q=\sum_{i=1}^{2k}P_i^2\). Define
\[
 F=\sum_{i=1}^{k}P_i^2-\sum_{i=k+1}^{2k}P_i^2,
 \qquad C^4_{N,k}=\{x\in\R^N:F(x)=0\}.
\]
This particular Clifford realization requires
$\delta(2k-1)\mid N/2$ \cite[Section~3.5]{FKM81}.
Here $F=f\circ P$, where
$f(y)=\sum_{i\leq k}y_i^2-\sum_{i>k}y_i^2$, so the composition
is the pullback of a Simons cone. The unequal-target example
\cref{ex:lawson-pullback} extends this connection with the transfer
theorem of \cref{sec:quartic-application}.
At the comparison endpoint \(k=1\),
\(F=(P_1-P_2)(P_1+P_2)\) is reducible and harmonic;
\cref{prop:minimizing-irreducibility} excludes area minimization.
For the minimizing parameters below, the same proposition gives
irreducibility.
The chain rule gives the identities recorded in
\cite[Lemma~5.2.11]{RiedlerThesis25}:
\begin{equation}\label{eq:quartic-ambient-closure}
 \bigl(\Gamma(I_i,I_j)\bigr)_{I=(R,Q,F)}=
 \begin{pmatrix}
 4R&8Q&8F\\
 8Q&16RQ&16RF\\
 8F&16RF&16RQ
 \end{pmatrix},\qquad
 \Delta I=(2N,16kR,0).
\end{equation}
The determinant is \(1024R(R^2-Q)(Q^2-F^2)\), so the spherical
quotient has dimension two wherever \(0<Q<R^2\) and \(|F|<Q\).
For a concrete example, take the quadratic octonionic Hopf map
\(\widetilde P=(P_1,\ldots,P_9):\R^{16}\to\R^9\), normalized by
\(|\widetilde P|^2=R^2\), and retain its first eight components.
Then \(k=4\) and \(Q=R^2-P_9^2\). Choose an image with
\(P_9\neq0\) and nonzero components in both four-dimensional blocks.
This gives \(0<Q<R^2\) and \(|F|<Q\). Taking two copies of this
eight-matrix Clifford system gives \(N=32\), \(k=4\), in the
minimizing range below; the same point, placed in one copy, retains
the strict inequalities.
When \(Q\not\equiv R^2\), these generators correspond to
\eqref{eq:orthogonal-quartic-data} via
\[
 (E_0,E_1,E_2)=\left(R^2-Q,\frac{Q+F}{2},\frac{Q-F}{2}\right),
 \qquad (r_0,r_1,r_2)=\left(\frac N2-2k+1,k,k\right).
\]
For $N\ge32$, $k\ge4$, normalize $T=Q/4$, $F_{\rm gen}=F/8$.
Lemma~\ref{lem:quartic-table-geometry} gives regular minimality,
the singular set $\{Q=0\}$, and full support. If $k=4$, then
$r_0\ge9$ and \cref{thm:quartic-table-existence} gives $G=RQF$;
if $k\ge5$, \cref{prop:balanced-quartic-multiplier} gives $G=QF$.
Thus the boundary current is strictly area minimizing in this range,
with the Hardy bound \eqref{eq:bihomogeneous-Hardy}. Positive
constant rescalings do not affect normalized gradients.

The defining polynomial itself satisfies
\[
 |\nabla F|^2=16RQ,\qquad \Delta F=0,\qquad
 \cL(F)=-64F(Q+2R^2).
\]
The last identity follows from
$\cL(F)=-\Gamma(F,\Gamma(F,F))/2$ and
\eqref{eq:quartic-ambient-closure}.

\begin{example}[Riedler's subcalibration potential]
\label{ex:Riedler-correct-power}
For \(N\geq32\) and \(k\geq4\), Riedler's thesis uses
\cite[Section~5.2.2, Lemmas~5.2.12--5.2.13]{RiedlerThesis25}
\begin{equation}\label{eq:Riedler-correct-potential}
 G=RQ^{3/4}F=16^{-3/4}R^{1/4}S^{3/4}F,
 \qquad S=|\nabla F|^2=16RQ.
\end{equation}
This potential has degree nine. The parameters in
\cref{prop:harmonic-power-Jacobi} are
\((e,\sigma,\gamma)=(4,1/2,3/4)\).
Using \eqref{eq:quartic-ambient-closure}, we obtain the boundary Jacobi coefficient
\begin{equation}\label{eq:Riedler-positive-boundary}
 \beta_G=(2N-41)+(12k-47)\frac{R^2}{Q}>0.
\end{equation}
Since \(Q\leq R^2\), \cref{thm:ground-state-gap} gives
\[
 Q_C(u)\geq
 \left[\left(9-\frac{N-1}{2}\right)^2+2N+12k-88\right]
 \int_Cr^{-2}u^2.
\]
This lower bound exceeds the general quartic bound
\eqref{eq:bihomogeneous-Hardy}, with $h=2k$, by
\[
 \left[\frac{(N-19)^2}{4}+2N+12k-88\right]
 -\left[\frac{(N-17)^2}{4}+16k-72\right]
 =N+2-4k=2r_0>0.
\]
Thus Riedler's degree-nine potential gives the stronger Hardy bound.
\end{example}

\begin{corollary}[Instability of Riedler's sixteen-dimensional quartic]
\label{thm:Riedler-unstable}
The cone $C^4_{16,4}$ is unstable on its regular part and is not area
minimizing. More precisely, there are compactly supported variations on the regular
part with scale-invariant quotients of the second variation tending to
$-15/4$.
\end{corollary}

\begin{proof}
Write $F_R$ for the defining quartic of $C^4_{16,4}$ and let
$A_1,\ldots,A_8$ be its symmetric Clifford matrices. The volume
matrix $A_9=A_1\cdots A_8$ satisfies
\[
 A_9^T=A_9,\qquad A_9^2=I,\qquad A_9A_i=-A_iA_9.
\]
Its two eigenspaces have dimension eight, since $A_1$ exchanges
them isometrically. Choosing compatible orthonormal bases gives
coordinates $(x,y)\in\R^8\oplus\R^8$ in which
\[
 A_{i+1}=\begin{pmatrix}0&B_i\\B_i^T&0\end{pmatrix}
 \quad(0\leq i\leq7),\qquad B_0=I.
\]
The matrices $B_1,\ldots,B_7$ are skew-symmetric and satisfy
$B_iB_j+B_jB_i=-2\delta_{ij}I$. Thus the eight vectors $B_i y$
$(0\leq i\leq7)$ are pairwise orthogonal, all of length $|y|$, and
\begin{gather*}
 \sum_{i=0}^7\langle x,B_i y\rangle^2=|x|^2|y|^2,
 \qquad Q=4|x|^2|y|^2,\\
 \frac{F_R}{8}
 =\sum_{i=0}^3\langle x,B_i y\rangle^2
       -\frac12|x|^2|y|^2=F_{\rm split}.
\end{gather*}
Consequently $C^4_{16,4}$ is orthogonally congruent to a splitting
cone \eqref{bh:centered} with $(\ell,m)=(8,4)$.
Its Clifford representation is the balanced one: the symmetric
involution $B_1B_2B_3$ anticommutes with $B_4$, so its two
eigenspaces have dimension four.

In the original variables, $\Lambda=C^4_{16,4}\cap\Sph^{15}$ and
$q=Q|_\Lambda=4t$, so \eqref{bh:link-table} reads
\begin{equation}\label{eq:Riedler-16-link}
 |A_\Lambda|^2=\frac{24}{q}+14,\qquad
 |\nabla_\Lambda q|^2=16q(1-q),\qquad
 \Delta_\Lambda q=56-64q.
\end{equation}
The $\ell=8$ argument above transfers under this orthogonal
identification, and \eqref{bh:negative-margin} gives the asserted
quotients and instability.
\end{proof}

This is the case left undecided in
\cite[Remark~5.2.6(3)]{RiedlerThesis25}. Together with
Theorem~5.2.5 and the exclusions in Remark~5.2.6(1)--(2)
of that thesis, it gives
\[
 C^4_{N,k}\text{ is area minimizing}
 \quad\Longleftrightarrow\quad N\geq32\text{ and }k\geq4
\]
among the admissible Clifford--Simons quartics.



\subsection{Further quartic realizations}

The remaining realizations in \cref{tab:quartic-weight-realizations}
have different singular geometry. Their subcalibrations and Hardy
bounds follow from the same quartic theorem.

\subsubsection{Two infinite families with real rank-one singular sets}
\label{sec:orthq-families}

An unused block of weight two gives two explicit nonharmonic
families in $\R^{8k}$. Their singular set is the real rank-one
matrix cone, and their area-minimizing range has the sharp threshold
$k\ge3$.

Identify $\R^{4k}$ with $\mathbb H^k$ and let $J_1,J_2,J_3$ be
componentwise right multiplication by the quaternionic units
$\mathbf i,\mathbf j,\mathbf k$. For $x,y\in\mathbb H^k$, $k\ge2$, put
$R=|x|^2+|y|^2$ and $W=|x|^2|y|^2-\langle x,y\rangle^2$, and define
\begin{align}
 F_{k,2}&=(4k-3)
  (\langle x\mathbf i,y\rangle^2+\langle x\mathbf j,y\rangle^2)-W,
  \label{eq:orthq-two-families}\\
 F_{k,3}&=(4k-3)
  (\langle x\mathbf i,y\rangle^2+\langle x\mathbf j,y\rangle^2
                         +\langle x\mathbf k,y\rangle^2)-2W.
  \nonumber
\end{align}
Here the displayed pairings are real Euclidean inner products.
For $h=\sum_\alpha\bar x_\alpha y_\alpha$, one has
$\langle x\mathbf i,y\rangle=\Im_{\mathbf i}h$ and similarly for
the other units. Thus the second family is also the Gram composition
\begin{equation}\label{eq:orthq-quaternionic-Gram}
 F_{k,3}=(4k-3)|h|^2-(4k-5)(\Re h)^2-2|x|^2|y|^2.
\end{equation}

\begin{theorem}[Two families and their sharp threshold]
\label{thm:orthq-infinite-families}
For $k\ge2$ and $m=2,3$, the polynomial $F_{k,m}$ defines an
irreducible minimal hypercone $C_{k,m}\subset\R^{8k}$ with the
algebra and singular geometry in \cref{prop:orthq-construction}.
The cone is area minimizing if and only if $k\ge3$.
For $k\ge3$, the degree-eight potential
\begin{equation}\label{eq:orthq-potential}
 G_{k,m}=WF_{k,m}
\end{equation}
subcalibrates the full boundary current, which is strictly area
minimizing. The regular cone satisfies
\begin{equation}\label{eq:orthq-Hardy}
 Q_{C_{k,m}}(\varphi)\ge
 \left(16k^2-36k-\frac{31}{4}\right)
 \int_{C_{k,m}}r^{-2}\varphi^2.
\end{equation}
At $k=2$ both cones are unstable on their regular parts.
\end{theorem}

Both sequences arise from the following construction, which gives
their algebra, irreducibility, and transverse singular geometry.

Let $J_1,\ldots,J_m$ be skew-symmetric orthogonal endomorphisms of
$\R^\ell$ satisfying
\[
 J_iJ_j+J_jJ_i=-2\delta_{ij}\Id,
 \qquad m\ge2,\qquad\ell-m-1\ge2.
\]
For $x,y\in\R^\ell$, put
\begin{gather}
 R=|x|^2+|y|^2,\qquad
 W=|x|^2|y|^2-\langle x,y\rangle^2,\qquad
 Q=\sum_{i=1}^m\langle J_ix,y\rangle^2,
 \label{eq:orthq-data}\\
 \begin{gathered}
 a=\ell-m-2,\qquad b=m-1,\qquad d=a-b=\ell-2m-1,\\
 F=(\ell-3)Q-bW.
 \end{gathered}
 \label{eq:orthq-polynomial}
\end{gather}
The vectors $x,J_1x,\ldots,J_mx$ are mutually orthogonal and have
length $|x|$, so $0\le Q\le W\le R^2/4$.

\begin{proposition}[Construction and exact algebra]
\label{prop:orthq-construction}
The polynomial $F$ is real irreducible and nonharmonic, has nonempty
regular zero set, and has essential space $\R^{2\ell}$.
Its regular zero set is minimal, and
\begin{equation}\label{eq:orthq-minimality}
 |\nabla F|^2=4R(dF+abW),\qquad
 \cL(F)=-8F\bigl(2dF+2abW+abR^2\bigr).
\end{equation}
Moreover,
\[
 \cA_R(F)=\R[R,W,F]
\]
is a polynomial algebra on three generators of degrees $(2,4,4)$.
Its critical set is the geometric singular set of $C=\{F=0\}$:
\begin{equation}\label{eq:orthq-singular}
 \Sing C=\{W=0\}
 =\{(x,y):\operatorname{rank}_{\R}[x\ y]\le1\}.
\end{equation}
This set has dimension $\ell+1$ and is smooth away from the origin.
At each nonzero singular point the transverse tangent cone is the
Lawson quadratic cone of signature $(m,\ell-m-1)$.
\end{proposition}

\begin{proof}
The symmetric Clifford system associated with
$\Id,J_1,\ldots,J_m$ has quadratic coordinates
$|x|^2-|y|^2$, $2\langle x,y\rangle$, and
$2\langle J_ix,y\rangle$; see \cite[Section~3]{FKM81}.
The blocks
\begin{equation}\label{eq:orthq-blocks}
 E_0=R^2-4W,\qquad E_1=4Q,\qquad E_2=4(W-Q)
\end{equation}
satisfy \eqref{eq:orthogonal-quartic-data} with weights
$(2,m,\ell-m-1)$ by \cref{lem:Clifford-block-quartics}. In the normalization of
\eqref{eq:quartic-general-generators}, $T=W$ and
$F_{\rm gen}=F/(\ell-3)$. Thus \eqref{bh:reduced-table} gives
\begin{equation}\label{eq:orthq-table}
 \begin{gathered}
 (\Delta R,\Delta W,\Delta F)
 =(4\ell,2(\ell-1)R,2dR),\\
 (\Gamma(I_i,I_j))_{I=(R,W,F)}=
 \begin{pmatrix}
 4R&8W&8F\\8W&4RW&4RF\\8F&4RF&4R(dF+abW)
 \end{pmatrix}.
 \end{gathered}
\end{equation}
Equation \eqref{eq:orthq-minimality} follows from
\eqref{bh:minimality}. The smallest-algebra and independence assertions
follow from the argument after \eqref{bh:reduced-table}.
In this case $R$ is already obtained from $F$ by one Laplacian:
$\ell$ is even because $J_1^2=-\Id$, whereas $2m+1$ is odd, so
$d\ne0$ and $R=\Delta F/(2d)$. In particular $F$ is nonharmonic.

For the remaining assertions, fix $x\ne0$ and decompose
\[
 \R^\ell=\R x\oplus U_x^+\oplus U_x^-,\qquad
 U_x^+=\operatorname{span}\{J_1x,\ldots,J_mx\}.
\]
For $y=y_0+y_++y_-$, the exact fibre formulas are
\begin{equation}\label{eq:orthq-fibre}
 W=|x|^2(|y_+|^2+|y_-|^2),\qquad
 F=|x|^2(a|y_+|^2-b|y_-|^2).
\end{equation}
They exhibit regular zeros and both phases. Writing
$F=y^{\mathsf T}M(x)y$, the matrix $M(x)$ has rank $\ell-1\ge4$
for every real $x\ne0$, with kernel $\R x$. It has the same rank
over $\R(x)$, so this quadratic in $y$ cannot factor into linear
forms over that field. Its coefficients, which are homogeneous
quadratics in $x$, have no nonconstant common divisor. A linear
common divisor would give $M(x)=0$ at a nonzero point of a real
hyperplane. A quadratic common divisor would give $M(x)=p(x)B$
for a constant matrix $B$, and $M(x)x=0$ would force $B=0$.
Both contradict the rank assertion. Gauss's lemma therefore proves
real irreducibility.

A constant annihilating direction $(v_1,v_2)$ would give two
vanishing derivatives of distinct bidegrees $(1,2)$ and $(2,1)$.
The second implies $M(x)v_2=0$ for every $x\ne0$, hence $v_2=0$;
interchanging $x,y$ gives $v_1=0$. Thus the ambient space is essential.
Finally $dF+abW=a^2Q+b^2(W-Q)$, so \eqref{eq:orthq-minimality}
identifies the critical set with $W=0$. Equality in real
Cauchy--Schwarz gives \eqref{eq:orthq-singular}. The charts $y=tx$
with $x\ne0$, and their counterparts with $y\ne0$, give its smoothness
and dimension. Smooth local frames in \eqref{eq:orthq-fibre} give
the transverse quadratic cone. At $(x,tx)$ an isometric normal
coordinate is $w=\Pi_{x^\perp}(\delta y-t\delta x)/\sqrt{1+t^2}$; the
nonzero Hessian eigenvalues are $2aR$ and $-2bR$, with multiplicities
$m$ and $\ell-m-1$. Their sum is at least four, so the zero set
cannot be smooth there. This also proves that these critical points
are geometric singularities, as in \cref{lem:quartic-table-geometry}.
\end{proof}

\begin{proof}[Proof of \cref{thm:orthq-infinite-families}]
The weights are $(2,m,4k-m-1)$, and their active sum is $h=4k-1$.
For $k\ge3$, this is at least eleven. The normalization
$T=W$, $F_{\rm gen}=F_{k,m}/(4k-3)$ therefore gives
\eqref{eq:orthq-potential} from \eqref{bh:potentials}, with $W>0$
off the rank-one singular set. The strict-minimality and full-support
conclusions follow from \cref{thm:quartic-table-existence}. In this normalization,
\begin{equation}\label{eq:orthq-potential-defect}
 \begin{split}
 \cL(WF)=8WF\bigl\{&[3(\ell-1)R^2-12W]F^2\\
 &+d[(\ell-5)R^2-4W]WF
 +ab[(\ell-9)R^2-4W]W^2\bigr\},
 \end{split}
\end{equation}
which is \eqref{bh:G8-defect} after the same constant rescaling.
The Hardy coefficient in \eqref{eq:bihomogeneous-Hardy}, with
$N=8k$ and $h=4k-1$, is exactly
\[
 \frac{(8k-17)^2}{4}+8(4k-1)-72
 =16k^2-36k-\frac{31}{4}>0.
\]

At $k=2$, the transverse tangent cones from
\cref{prop:orthq-construction} have signatures $(2,5)$ and $(3,4)$
in $\R^7$. Each is a six-dimensional Lawson cone whose link has
$|A|^2=5$, so its radial Hardy margin is
$(6-2)^2/4-5=-1$. \Cref{lem:tangent-instability-transfer},
applied to the flat $\R^9$
factor and the smooth local model \eqref{eq:orthq-fibre}, transfers
a compactly supported negative variation to the regular cone.
Thus both $k=2$ cones are unstable and not area minimizing.
\end{proof}

\begin{proposition}[Distinction from the preceding families]
\label{prop:orthq-distinct}
For $k\ge3$, the cones in \eqref{eq:orthq-two-families} are not
orthogonally congruent to the standard OT--FKM splitting cones or
to pure full Clifford pullbacks of nondegenerate Lawson quadratic
cones. For each fixed $k$, the two cones are not congruent to one another.
\end{proposition}

\begin{proof}
The splitting cones in $\R^{2\ell}$ have singular set the two
coordinate $\ell$-planes, whereas \eqref{eq:orthq-singular} has
dimension $\ell+1$. For a full Clifford quadratic map
$P:\R^{2\ell}\to\R^p$, the identity
$DP\,DP^{\mathsf T}=4R\Id_p$ shows that the nonvertex singular
set of a nondegenerate Lawson pullback is $P^{-1}(0)$, of codimension
$p$ if nonempty. For the balanced target this is also recorded in
\cite[Remark~5.2.6(1)]{RiedlerThesis25}. Congruence would force $p=\ell-1$.
The Clifford module restriction
\cite[Section~3.5]{FKM81}, equivalently
\cite[Table~4.1]{RiedlerThesis25}, would then require \(\delta(\ell-2)=2^{2k-1}\mid4k=\ell\),
which is impossible for $k\ge3$ because $2^{2k-1}>4k$.
Finally the transverse signatures $(2,4k-3)$ and $(3,4k-4)$
are distinct even after interchanging their two entries, and hence
distinguish the two sequences.
\end{proof}

\begin{remark}[The weights $(2,5,10)$]
\label{ex:quartic-2510}
Taking $\ell=16$ and $m=5$ in \cref{prop:orthq-construction}
gives the weights $(2,5,10)$ and $F=13Q-4W$ on $\R^{32}$.
The required Clifford system exists since $\delta(6)=8\mid16$.
Here $h=15$, so \cref{thm:quartic-table-existence} gives the
potential $WF$, strict area minimization, and strict stability.
\end{remark}

\subsubsection{A complex rank-one singular set: weights \texorpdfstring{$(3,7,7)$}{(3,7,7)}}
\label{sec:weighted-Clifford-quartic}

The balanced criterion \cref{prop:balanced-quartic-multiplier} gives
a harmonic example with a complex rank-one singular set.
Let $E_i=L_{e_i}$, $1\leq i\leq7$, be left multiplication by an
orthonormal imaginary basis of the octonions.  Alternativity gives
$E_i^T=-E_i$ and $E_iE_j+E_jE_i=-2\delta_{ij}\Id_8$.
On $\R^{16}=\R^8\oplus\R^8$ put
\begin{equation}\label{eq:weighted-Clifford-matrices}
 J=\begin{pmatrix}0&\Id_8\\-\Id_8&0\end{pmatrix},\qquad
 B_i=\begin{pmatrix}E_i&0\\0&-E_i\end{pmatrix}.
\end{equation}
These eight skew-symmetric matrices satisfy the negative Clifford
relations.  For $x,y\in\R^{16}$ define
\begin{align}
 R&=|x|^2+|y|^2,\nonumber\\
 W&=|x|^2|y|^2-\langle x,y\rangle^2-\langle Jx,y\rangle^2,
 \label{eq:weighted-Clifford-definition}\\
 F&=2\sum_{i=1}^7\langle B_ix,y\rangle^2-W.\nonumber
\end{align}

\begin{proposition}[The complex rank-one example]
\label{thm:weighted-Clifford-quartic}
The polynomial $F$ in \eqref{eq:weighted-Clifford-definition} is
real irreducible and harmonic. Its cone $C=\{F=0\}\subset\R^{32}$ has
geometric singular set
\begin{equation}\label{eq:weighted-Clifford-singular-set}
 \Sing C=\Sigma=\{W=0\}
 \cong\{(z,w)\in\mathbb C^8\times\mathbb C^8:
             \operatorname{rank}_{\mathbb C}[z\ w]\leq1\}.
\end{equation}
This set has real dimension $18$ and is smooth away from the origin.
At each nonzero singular point the tangent cone is orthogonally
congruent to $\R^{18}\times\{|\xi|^2=|\eta|^2\}$, with
$\xi,\eta\in\R^7$. The potential $G=WF$ subcalibrates the full
boundary current, which is strictly area minimizing. Its smallest
Laplacian algebra is $\R[R,W,F]$, and \eqref{eq:bihomogeneous-Hardy} gives
\begin{equation}\label{cor:weighted-Clifford-Hardy}
 Q_C(\varphi)\geq\frac{385}{4}\int_C r^{-2}\varphi^2,
 \qquad \varphi\in C_c^\infty(C_{\reg}\setminus\{0\}).
\end{equation}
\end{proposition}

\begin{proof}
The Clifford system associated with $\Id_{16},J,B_1,\ldots,B_7$
has quadratic coordinates $p_0=|x|^2-|y|^2$ and
$p_A=2\langle Ax,y\rangle$. Apply \cref{lem:Clifford-block-quartics} with
\begin{align}
 E_0&=p_0^2+p_{\Id}^2+p_J^2=R^2-4W,\nonumber\\
 E_+&=\sum_{i=1}^7p_{B_i}^2=2(W+F),\label{eq:weighted-Clifford-blocks}\\
 E_-&=R^2-E_0-E_+=2(W-F).\nonumber
\end{align}
Here $\ell=16$, $m=9$, and the weights are $(3,7,7)$, so
\cref{prop:balanced-quartic-multiplier} gives harmonicity, the smallest
algebra, and the subcalibration. Irreducibility follows from
\cref{prop:minimizing-irreducibility}; the Hardy coefficient uses
$(N,h)=(32,14)$.

Identify $(\R^{16},J)$ with $\mathbb C^8$.  The expression for
$W$ is the complex Gram determinant of the two columns $x,y$.
Thus \cref{lem:quartic-table-geometry} and equality in Cauchy--Schwarz
give \eqref{eq:weighted-Clifford-singular-set}.
The rank-one matrices form a smooth complex manifold of dimension
$8+2-1=9$, so $\dim_{\R}\Sigma=18$ and
$\Sigma\setminus\{0\}$ is smooth.

For $x\ne0$, the vectors $x,Jx,B_1x,\ldots,B_7x$ are orthogonal
with length $|x|$. Decompose $y=y_0+y_++y_-$ along
$K_x=\operatorname{span}\{x,Jx\}$,
$U_x^+=\operatorname{span}\{B_1x,\ldots,B_7x\}$, and
$U_x^-=(K_x\oplus U_x^+)^\perp$. The last two spaces have dimension seven, and
\begin{equation}\label{eq:weighted-Clifford-transverse-form}
 W=|x|^2(|y_+|^2+|y_-|^2),\qquad
 F=|x|^2(|y_+|^2-|y_-|^2).
\end{equation}
Local frames give transverse signature $(7,7)$; interchanging $x,y$
covers $\Sigma\setminus\{0\}$. The Hessian formulas in
\cref{lem:quartic-table-geometry}, with $T=W$ and $F_{\rm gen}=F/2$,
give eigenvalues $\pm2R(p)$, each sevenfold, and kernel dimension
eighteen; \cref{sec:tangent-instability-transfer} gives the asserted
orthogonal tangent-cone model.
\end{proof}

The singular dimension $18$ distinguishes this cone from the
splitting and Hopf pairing cones in $\R^{32}$, whose singular sets
have dimension $16$. For a pure full Clifford pullback of a
nondegenerate Lawson cone, \eqref{eq:Clifford-pullback-identities}
makes the nonvertex singular set the zero fiber of the Clifford map.
Congruence would therefore require fourteen quadratic coordinates.
Such a system on $\R^{32}$ would require $\delta(13)=128\mid16$,
contrary to the module restriction \cite[Section~3.5]{FKM81}.
In particular, the example is distinct from Riedler's family.


\subsubsection{An unequal Lawson pullback}

\begin{example}[An unequal quadratic target]
\label{ex:lawson-pullback}
A Lawson cone in $\R^5\oplus\R^6$ pulls back to a quartic in
$\R^{64}$, illustrating \cref{prop:pullback-subcalibration}.
On the target, put
\[
 T=|u|^2+|v|^2,\qquad f=5|u|^2-4|v|^2,\qquad H=Tf.
\]
The target potential satisfies \eqref{eq:pullback-target-sign} with
$\ell=4$ and $\kappa=1/2$, strictly off $f=0$; the calculation is
recorded in \cref{app:application-details}.

Choose eleven Clifford matrices on \(\R^{64}\), and set
\(F=f\circ P\), \(Q=T\circ P=|P|^2\), and
\(R=|x|^2\) on the source. Thus \(G=H\circ P=QF\).
The source blocks
\[
 E_0=R^2-Q,\qquad E_1=\sum_{i=1}^{5}P_i^2,\qquad
 E_2=\sum_{i=6}^{11}P_i^2
\]
have weights $(22,5,6)$. Such matrices exist by
\cite[Table~4.1]{RiedlerThesis25}. With $T_{\rm gen}=Q/4$,
$F_{\rm gen}=F/36$, \cref{thm:quartic-table-existence} gives the
same potential $QF$, critical set $P^{-1}(0)$, full support, and
strict area minimality. Here $Q>0$ off $P^{-1}(0)$, so the
multiplier has the required positive domain. The Hardy coefficient
in \eqref{eq:bihomogeneous-Hardy} is $2273/4$.
\end{example}


\subsection{An example from complex Gram coordinates}
\label{sec:gram-composition-constructions}

The following sextic is obtained from the nine quadratic coordinates
of a complex Hermitian Gram matrix. Its profile uses the imaginary
part of a cyclic product of three off-diagonal entries. It shows
that the composition criterion also produces higher-degree minimal
regular zero sets.

\begin{example}[A complex Gram sextic]
\label{thm:complex-gram-cycle-sextic}
For every $k\ge2$, let $x_1,x_2,x_3\in\mathbb C^k$ and
$h_{ij}=x_i^*x_j$.  Then
\begin{equation}\label{eq:complex-gram-cycle-sextic}
 F=\operatorname{Im}(h_{12}h_{23}h_{31})
\end{equation}
is a real harmonic homogeneous polynomial of degree six on $\R^{6k}$
with nonempty minimal regular zero set.

More explicitly, set $r_i=h_{ii}$, $R=\sum_i r_i$,
\[
 a_1=|h_{23}|^2,\qquad a_2=|h_{31}|^2,\qquad
 a_3=|h_{12}|^2,\qquad
 t=\operatorname{Re}(h_{12}h_{23}h_{31}).
\]
With $\{i,j,l\}=\{1,2,3\}$ in each summand, define
\begin{align}
 S&=\sum_{i=1}^3 r_i a_i(a_j+a_l)-2t\sum_{i=1}^3a_i,
 \label{eq:complex-gram-cycle-gradient}\\
 B&=\sum_{i=1}^3a_i(r_i+r_j)(r_i+r_l)
       -\sum_{i=1}^3a_i^2-2Rt.
 \label{eq:complex-gram-cycle-weight}
\end{align}
Then
\begin{equation}\label{eq:complex-gram-cycle-minimality}
 \Delta F=0,\qquad |\nabla F|^2=S,\qquad
 \Gamma(F,S)=4BF,\qquad \cL(F)=-2BF.
\end{equation}
\end{example}

\begin{proof}
Use the complex bilinear extension of $\Gamma$ in
\eqref{eq:complex-Gram-differential-table}.  For
$Z=h_{12}h_{23}h_{31}$ and $A=a_1+a_2+a_3$, the product rule gives
\[
 \Gamma(Z,Z)=4AZ,\qquad
 \Gamma(Z,\overline Z)=2\sum_i r_i a_i(a_j+a_l).
\]
Taking real and imaginary parts proves
$|\nabla F|^2=S$.  The same table gives
\begin{equation}\label{eq:complex-gram-cycle-pairings}
 \Gamma(F,r_i)=4F,\quad
 \Gamma(F,a_i)=2(R-r_i)F,\quad
 \Gamma(F,t)=2AF.
\end{equation}
Applying these three identities to the expression for $S$ gives
$\Gamma(F,S)=4BF$. Indeed, writing
$D=\sum_i r_i a_i$ and $E=\sum_i r_i a_i^2$, one has
$S=AD-E-2tA$. After the $A^2$ terms cancel, the product rule gives
\[
 \Gamma(F,S)=4F\left[
 RD+\frac A2\sum_i r_i(R-r_i)-\sum_i a_i^2
 -\sum_i r_i a_i(R-r_i)-2Rt\right].
\]
The coefficient of $a_i$ in the brackets is
$r_i^2+r_1r_2+r_2r_3+r_3r_1=(r_i+r_j)(r_i+r_l)$,
so the brackets equal $B$. This is a polynomial identity, valid
also when the Gram entries satisfy rank relations, as for $k=2$.
For harmonicity, with $x_2,x_3$ fixed, $F$ is the Hermitian
quadratic form in $x_1$ with matrix
\[
 \frac{h_{23}x_2x_3^*-\overline{h_{23}}x_3x_2^*}{2i}.
\]
Its trace is zero, so its block Laplacian vanishes.  The same
argument applies to the other two blocks.  The point
$x_1=x_3=e_1$, $x_2=e_1+e_2$ has $F=0$ and $S=2$, proving
nonempty regularity; \eqref{eq:complex-gram-cycle-minimality}
then proves minimality.

\end{proof}



\subsection{Higher bidegrees and transverse instability}

The two-block equations also admit higher bidegrees. In the following
octonionic examples, an unstable transverse tangent cone excludes
area minimization.

Identify $\R^8$ with $\mathbb O$ and set
\[
 \mu(x,y)=x\overline y,\qquad x,y\in\mathbb O.
\]
Then $\norm{\mu}^2=T$ and
\begin{equation}\label{bh:mu}
 D\mu\,D\mu^{\mathsf T}=R\Id_8,\qquad
 \Delta\mu=0,\qquad D\mu(\nabla R)=4\mu.
\end{equation}
These follow because multiplication by a fixed nonzero octonion is
a linear similarity, with its squared norm as dilation squared;
see \cite[Chapter~1, Lemma~1.3(1) and Appendix, pp.~15--16 and 21--22]{Shapiro2000}
for the norm-multiplicative algebra and similarity identities.

We use T.~Wang's octonionic pullback construction
\cite[Section~3.4]{WangTeng26}. For $g=3,4,6$, choose the normalized
Cartan--M\"unzner polynomial $\Phi_g$ on $\R^8$ with multiplicities
respectively $(2,2)$, $(1,2)$, and $(1,1)$. Set
\[
 P_g=\Phi_g\circ\mu,\qquad F_g=P_g-t_gT^{g/2},\qquad
 (\kappa_g,t_g)=(0,0),(8,1/3),(0,0),
\]
in the same order, where $\Delta\Phi_g=\kappa_g|z|^{g-2}$ and
terms with zero coefficient are omitted. The value $t_g$ is the
minimal level from \cref{sec:low-rank}. These polynomials have
bidegree $(g,g)$.

\begin{proposition}\label{bh:octonionic}
Each $\R[u,v,P_g]$ is a graded Laplacian algebra with three
algebraically independent generators. Each $C_g=\{F_g=0\}\subset\R^{16}$
is a minimal hypercone, with singular set equal to the two coordinate
$\R^8$ subspaces. For $g=3,4,6$, these hypercones are unstable.
\end{proposition}
\begin{proof}
The closure identities are
\begin{gather}\label{bh:higher-table}
 \Delta u=\Delta v=16,\quad\Gamma(u,v)=0,\quad
 \Gamma(u,u)=4u,\quad\Gamma(v,v)=4v,\\
 \Gamma(u,P_g)=\Gamma(v,P_g)=2gP_g,\quad
 \Gamma(P_g,P_g)=g^2RT^{g-1},\nonumber\\
 \Delta P_g=\kappa_gRT^{g/2-1}.\nonumber
\end{gather}
The last expression is polynomial (zero for $g=3$).
The gradient determinant $16g^2R(T^g-P_g^2)$ is positive at
intermediate levels. Thus \cref{prop:MR-generator-test} gives
closure and \eqref{eq:background-quotient-Gram} gives independence
and spherical quotient dimension two.

Put $H_g(z)=\Phi_g(z)-t_g\norm{z}^g$, with the second term zero when
$t_g=0$. This is a homogeneous polynomial in every case under
consideration, and its zero set is the minimal isoparametric
hypercone in $\R^8$. A useful consequence of \eqref{bh:mu} is, for any
homogeneous $H$ of degree $h$,
\begin{equation}\label{bh:mu-pullback}
 \cL(H\circ\mu)=R^2(\cL(H))\circ\mu
                    -2h(H\norm{\nabla H}^2)\circ\mu.
\end{equation}
Indeed $\norm{\nabla(H\circ\mu)}^2=R(\norm{\nabla H}^2)\circ\mu$,
$\Delta(H\circ\mu)=R(\Delta H)\circ\mu$, and
$\Gamma(H\circ\mu,R)=4h(H\circ\mu)$. Substitution proves
\eqref{bh:mu-pullback}. Both terms vanish on $H_g\circ\mu=0$, proving
minimality. Off the coordinate subspaces $\mu\ne0$, the target
minimal level is regular, and \eqref{bh:mu} gives a nonzero source
gradient. Along either coordinate subspace that gradient vanishes.

At $(x_0,0)$ with $x_0\ne0$, the first nonzero term is
$H_g(x_0\overline y)$, so the tangent cone is a product of $\R^8$
with an orthogonal image of the target minimal cone. The target
link has dimension six and constant squared curvature $6(g-1)$.
Its classical Hardy margin is
\[
 D_g=\frac{25}{4}-6(g-1)=
 \begin{cases}-23/4,&g=3,\\-47/4,&g=4,\\-95/4,&g=6.
 \end{cases}
\]
Each is negative, so \cref{lem:tangent-instability-transfer}
proves instability of $C_g$. The target cone
is nonlinear, so these coordinate strata are genuine singularities.
\end{proof}

\begin{remark}[The quadratic target and the sign obstruction]
For a quadratic target with both multiplicities positive, let
$Q=\norm{\operatorname{pr}_W\mu}^2$, $m=\dim W$, and
$c=(m-1)/6$. This $Q$ obeys \eqref{bh:original-table} with $\ell=8$.
The curvature and cutoff proof above applies without a choice of
Clifford splitting, so the full $(2,2)$ hypercone $Q=cT$ is also
unstable. Thus all four bidegrees in this $\R^8\times\R^8$ construction
give unstable hypercones. The minimizing question for a codimension-two balance section is separate.
\end{remark}

